\documentclass[UTF8,11pt,fontset=none]{ctexart}
\usepackage[a4paper,margin=25mm,headheight=15pt]{geometry}
\usepackage{fontspec}
\setCJKmainfont[ItalicFont={Songti SC}]{Songti SC}
\setCJKsansfont{Heiti SC}
\setCJKmonofont{PingFang SC}
\usepackage{amsmath,amssymb,amsthm,mathtools,bm,mathrsfs}
\usepackage{booktabs,longtable,array,enumitem,microtype}
\usepackage{xcolor}
\usepackage[unicode,colorlinks=true,linkcolor=blue!45!black,citecolor=blue!45!black,urlcolor=blue!45!black]{hyperref}
\usepackage[numbers,sort&compress]{natbib}
\usepackage{fancyhdr}
\pagestyle{fancy}\fancyhf{}\fancyhead[L]{\small 分数次极大算子的尺度不变光晕}\fancyhead[R]{\small 唐浩天}\fancyfoot[C]{\thepage}
\theoremstyle{definition}
\newtheorem{theorem}{定理}[section]
\newtheorem{lemma}[theorem]{引理}
\newtheorem{proposition}[theorem]{命题}
\newtheorem{corollary}[theorem]{推论}
\newtheorem{conjecture}[theorem]{猜想}
\theoremstyle{definition}
\newtheorem{definition}[theorem]{定义}
\newtheorem{remark}[theorem]{注}
\newtheorem{example}[theorem]{例}
\newtheorem{question}[theorem]{问题}
\renewcommand{\proofname}{证明}
\numberwithin{equation}{section}
\setlist{nosep}
\allowdisplaybreaks[2]
\setlength{\emergencystretch}{2em}
\title{分数次极大算子的尺度不变光晕理论\\\large 双密度刚性、切向几何、带权刻画与多线性传递}
\author{唐浩天}
\date{2026年11月4日修订}
\begin{document}
\maketitle
\begin{abstract}
本文研究分数次极大算子在尺度不变阈值 $\alpha|E|^{\beta/n}$ 下的光晕。
严格见证不等式同时约束局部密度和全局捕获比例，并给出最佳代数对称差误差。
结合近重合几何，证明非中心方体和强矩形的端点增量为 $1-\alpha$ 阶，
中心方体为 $(1-\alpha)^n$ 阶。通过共同核质量修改、面层压缩及凸纤维对称平均，
进一步确定中心方体在全体可测输入上的精确首项，并给出紧分数阶区间上的统一下一阶余项。对单位方体，联合优化中心和平移尺度，得到显式切向
多面体及其体积 $2^n/[k!(2c)^{n-k}]$，其中 $c=\min\{\beta,n-\beta\}$、
$k=\lfloor2c\rfloor$。一维情形求得高阈值精确常数；外层权情形建立集中模量的双向
控制和 $A_\infty/A_\infty^\ast$ 刻画。多线性对角模型精确约化到线性参数 $\beta/m$，
一般输入则满足几何平均归一化下的双边夹逼。最后推导全体中心限制型高矩首项、分数阶
消失时的过渡下剖面，以及 Riesz、Wolff 势的最佳方体比较和直接应用。
一维闭式还给出全阶统一余项和临界窗口的四阶展开。对任意固定满维锚定凸体，进一步建立尺度与平移联合切向极限，并求得固定方体的加权首项。方体内部的尺度优化可精确有限化；二维固定模型由此得到四分支面积公式、全部临界窗口的过渡系数，以及带显式余项界的任意阶统一边界截断表示。
任意维固定方体还具有显式次首项，以及任意阶统一径向截断；其误差阶在一般参数类中不可提高。
在中心对称半阶模型中，切向体恰为极投影体，由Petty不等式确定椭球极值。
二维半阶由精确两点差体得到显式体积上包络；独立八尺度的交并预算与各向异性侧弧传递进一步证明，该包络与真实全体常数之间的亏损阶恰为 $(1-\alpha)^3$，并给出在 $0<1-\alpha\le2^{-14}$ 时的显式减项 $2^{-68}(1-\alpha)^3$，从而使全体三阶上极限严格小于 $10/3$。矩形下界给出下极限至少 $2$；精确三阶系数及极限存在性仍待研究。进一步对与正方形对称差为 $O((1-\alpha)^2)$ 的任意输入证明统一三阶上界，四阶余项阶数不可降低；偏心参数为 $\rho$ 的有间隙近矩形类还具有三阶上界 $2+2|\rho|(1-|\rho|)\le5/2$；有界角孔具有等号刚性和定性稳定性。有限四轴见证还迫使任意输入在测度上二阶接近矩形；方向变化的共同尺度八点容量具有解析边界与显式十字形最优输入。
\end{abstract}
\noindent\textbf{关键词：}分数次极大算子；尺度不变光晕；Tauberian 常数；切向多面体；外层权；多线性传递。

\clearpage
\tableofcontents
\clearpage
\section{引言}\label{sec:introduction}

几何极大算子的光晕描述一个集合在高平均阈值下的扩张。微分基的经典理论
\cite{deGuzman1975}与 Solyanik 的工作\cite{Solyanik1993}说明，阈值接近顶层时的
Tauberian 常数包含微分和覆盖几何的信息。设 $0<\beta<n$，$\mathcal B$ 是轴平行
非中心方体基 $\mathcal Q$、强矩形基 $\mathcal R$ 或中心方体基 $\mathcal Q_c$。
本文的对象是
\[
 M_{\beta,\mathcal B}f(x)=\sup_{B\in\mathcal B(x)}
 |B|^{\beta/n-1}\int_B|f|,
 \qquad
 C_{\beta,\mathcal B}(\alpha)=\sup_{0<|E|<\infty}
 \frac{|\{M_{\beta,\mathcal B}\mathbf1_E>\alpha|E|^{\beta/n}\}|}{|E|}.
\]
所有平均和阈值均以 Lebesgue 测度定义。方体的各边平行于坐标轴；中心基只有一个
共同边长，强矩形允许各边独立变化。

经典零阶光晕的 Solyanik 估计、局部 H\"older 正则性和外层权刻画分别见
\cite{HagelsteinParissis2014,HagelsteinParissis2015,HagelsteinParissis2016,HagelsteinParissis2018}。
分数次算子在欧氏伸缩下带有 $\beta$ 次齐次度，因此集合大小必须进入阈值。
固定的原始阈值会使全体集合的 Tauberian 商发散；保留分母 $|E|$ 时，幂型阈值的
伸缩不变性唯一决定指数 $\beta/n$。确定这一归一化之后，关键问题变为：顶层见证
能偏离资料集多少，见证中心又能在多大区域内变化？

已有相邻问题需要按对象区分。Aldaz\cite{Aldaz2021}研究任意测度下中心极大算子的
均匀弱型与交叠几何；本文实际使用的弱型界则在 Lebesgue 测度上直接证明。
分数次两权范数\cite{BanuelosOsekowski2025}、限制反向弱型\cite{Lerner2026}与
势算子的端点不可改进性\cite{MihulaPickSchikorra2026}提供相关背景，其归一化和
优化对象各有不同。卷积体的局部几何\cite{Haddad2025}及仿射等周理论
\cite{LutwakYangZhang2000}与固定凸体测试有关；从固定体的极限到对所有可测集的
统一最优首项，还需要形状稳定性和匹配上界。

本文的基本引理是双密度分解。记 $\theta=\beta/n$。若 $B$ 严格见证光晕中的一点，
并置 $s=|E\cap B|$、$a=s/|B|$、$b=s/|E|$，则
\[
 a^{1-\theta}b^\theta>\alpha,
 \qquad
 \frac{|B\triangle E|}{|E|}
 <\rho_\theta(\alpha):=
 \max\{1-\alpha^{1/\theta},\alpha^{-1/(1-\theta)}-1\}.
\]
这个误差是仅依赖见证不等式时的最佳值。它把分析问题转成近重合几何：任何两个
高水平见证必须接近同一集合，从而彼此接近。非中心基允许见证内的全部点作为
光晕点，中心基只允许见证的中心；由此得到不同的端点幂指数
\[
 C_{\beta,\mathcal Q}(\alpha)-1\asymp1-\alpha,
 \quad C_{\beta,\mathcal R}(\alpha)-1\asymp1-\alpha,
 \quad C_{\beta,\mathcal Q_c}(\alpha)\asymp(1-\alpha)^n.
\]

对固定单位方体 $Q_0$，进一步让边长和中心同时作一阶变化，可以消去尺度参数，
得到由有序坐标和定义的切向多面体 $S_{n,c}$。定理~\ref{tan:unit-cube-limit}证明
\[
 |H_{\beta,\mathcal Q_c}(1-t,Q_0)|
 =\frac{2^n}{k!(2c)^{n-k}}t^n+o(t^n).
\]
该固定输入极限产生全体中心常数的下界。在一维，区间相交和直径估计则给出真正的
全体常数闭式。外层权的集中模量由同一见证聚拢机制控制，反向以挖去小集合的
测试恢复集中条件。多线性部分利用非负单因子不超过顶层这一性质，得到精确对角
恒等式与一般输入夹逼，无须另建一套覆盖理论。

这些结论还给出归一化示性资料的高矩渐近。势论应用保留单向比较所能直接推出的
超水平集和集中限制；有关势估计、测度资料方程与分布函数方法的背景见
\cite{KilpelainenMaly1994,KuusiMingione2014,PhucVerbitsky2022,NguyenTran2021}，
Riesz 势弱型常数的相关研究见\cite{HuangTang2023}。

由一维闭式，命题~\ref{one:all-orders}给出任意阶展开及紧参数集上的显式余项；命题~\ref{one:transition}确定有限阈值转换曲线，并在临界窗口内给出四阶一致展开。

定理~\ref{cvx:full-tangent-limit}将固定尺度的卷积体下界推进为任意满维锚定凸体的联合切向极限，尺度参数由一个边界缺陷泛函消去。定理~\ref{wgt:theorem-fixed-cube-tangent}同时给出局部可积奇性权下的固定方体首项。

定理~\ref{fix:finite-scales}利用对数尺度上的分段凸性，把固定方体的内部尺度优化归约到有限接触尺度。二维进一步得到精确面积公式；例如固定单位正方形在 $\beta=1$ 时，充分高阈值下的光晕面积为 $2(1-\alpha)^2/\alpha$。

在 $K=-K$、$\beta=n/2$ 时，定理~\ref{half:projection-identity}精确消去全部尺度变化，切向体等于 $|K|\Pi^*K$。结合经典Petty投影不等式，可以在这一固定模型子类中确定椭球为最优资料体。

全文先统一定义并证明代数刚性，再建立近重合几何、方体切向极限和固定凸体推广。
随后处理一维常数、外层权和多线性模型，最后把端点光晕用于高矩、双参数过渡与
势函数比较。高维全体中心首项及近极值分类作为具体问题列于应用章末；它们与
经典 Halo 猜想的研究背景\cite{Hagelstein2025}保持相应的对象区分。

\section{预备知识与尺度归一化}
\label{sec:preliminaries}

本节固定 $n\geq1$ 与 $0<\beta<n$，并记
\[
 \theta:=\frac\beta n\in(0,1).
\]
设 $\mathcal Q$ 为所有轴平行方体构成的基，$\mathcal R$ 为所有轴平行矩形构成的基，$\mathcal Q_c$ 为以取值点为中心的轴平行方体构成的基。对其中任一基 $\mathcal B$，令 $\mathcal B(x)$ 表示可在 $x$ 使用的集合，并定义
\begin{equation}
 M_{\beta,\mathcal B}f(x):=
 \sup_{B\in\mathcal B(x)}|B|^{\theta-1}\int_B|f(y)|\,dy.
 \label{fnd:def-max}
\end{equation}
以下所有集合都取 Lebesgue 可测；$|\cdot|$ 表示 Lebesgue 测度。

\begin{lemma}[严格水平集的可测性]
\label{fnd:measurability}
若 $f\in L^1_{\mathrm{loc}}(\mathbb R^n)$ 与 $\lambda>0$，则对
$\mathcal B\in\{\mathcal Q,\mathcal R,\mathcal Q_c\}$，集合
\[
 \{x:M_{\beta,\mathcal B}f(x)>\lambda\}
\]
是开集。
\end{lemma}

\begin{proof}
先取 $\mathcal B=\mathcal Q_c$。对每个 $\ell>0$，令
$Q_\ell=[-\ell/2,\ell/2]^n$ 及
\[
 F_\ell(x):=\int_{x+Q_\ell}|f(y)|\,dy.
\]
函数 $F_\ell$ 连续：若 $x_j\to x$，则所有 $x_j+Q_\ell$ 与 $x+Q_\ell$ 落在同一有界方体内；在此方体上以 $|f|$ 支配，且平移后的方体指示函数在边界之外逐点收敛，故可由支配收敛定理得到结论。因此
\[
 \{M_{\beta,\mathcal Q_c}f>\lambda\}
 =\bigcup_{\ell>0}\{x:\ell^{\beta-n}F_\ell(x)>\lambda\}
\]
为开集。

再令 $\mathcal B=\mathcal Q$ 或 $\mathcal R$，并设 $x$ 有一个严格见证矩形
$B=\prod_{j=1}^n[a_j,b_j]$。对 $t\geq0$ 记
\[
 B_t:=\prod_{j=1}^n[a_j-t,b_j+t],\qquad
 \Phi(t):=|B_t|^{\theta-1}\int_{B_t}|f|.
\]
当 $t\downarrow0$ 时，$\Phi(t)\to\Phi(0)>\lambda$。于是某个 $t>0$ 满足 $\Phi(t)>\lambda$；并且 $x\in\operatorname{int}B_t$。因此 $x$ 的一个邻域中的每一点均可把同一个 $B_t$ 作为严格见证，所求水平集开。
\end{proof}

\begin{lemma}[示性函数的自然顶层]
\label{fnd:top-level}
若 $0<|E|<\infty$，则对任意 $x\in\mathbb R^n$ 及
$\mathcal B\in\{\mathcal Q,\mathcal R,\mathcal Q_c\}$，
\[
 M_{\beta,\mathcal B}\mathbf1_E(x)\leq |E|^\theta.
\]
\end{lemma}

\begin{proof}
对任一见证 $B$ 写 $s=|E\cap B|$。由于 $s\leq |B|$ 及 $s\leq |E|$，
\[
 s=s^{1-\theta}s^\theta
 \leq |B|^{1-\theta}|E|^\theta.
\]
乘以 $|B|^{\theta-1}$ 后取上确界即可。
\end{proof}

\begin{definition}[尺度不变光晕与 Tauberian 常数]
\label{fnd:def-halo}
对 $0<\alpha<1$，令
\[
 H_{\beta,\mathcal B}(\alpha,E):=
 \{x:M_{\beta,\mathcal B}\mathbf1_E(x)>\alpha|E|^\theta\},
 \qquad
 C_{\beta,\mathcal B}(\alpha):=
 \sup_{0<|E|<\infty}\frac{|H_{\beta,\mathcal B}(\alpha,E)|}{|E|}.
\]
\end{definition}

\begin{remark}
\label{fnd:center-density}
分数阶情形中，$x$ 是 $E$ 的 Lebesgue 密度点并不推出 $x\in H_{\beta,\mathcal Q_c}(\alpha,E)$：半径趋于零的中心方体平均量级为 $|Q|^\theta$，其极限为零。因此中心常数的端点行为不应和零阶极大算子混同。
\end{remark}

\begin{proposition}[伸缩律与唯一归一化]
\label{fnd:scaling}
对 $r>0$，有
\[
 M_{\beta,\mathcal B}\mathbf1_{rE}(rx)
 =r^\beta M_{\beta,\mathcal B}\mathbf1_E(x),
 \qquad
 H_{\beta,\mathcal B}(\alpha,rE)=rH_{\beta,\mathcal B}(\alpha,E).
\]
此外，若保持分母为 $|E|$，把阈值写成 $\alpha|E|^\sigma$，并要求变换
$(E,x)\mapsto(rE,rx)$ 对任意 $r>0$ 保持同一个参数 $\alpha$，则必且只可取
$\sigma=\beta/n$。
\end{proposition}

\begin{proof}
每个 $rx$ 处的见证均唯一写成 $rB$。作变量代换 $y=rz$，则
\[
 |rB|^{\theta-1}\int_{rB}\mathbf1_{rE}(y)\,dy
 =r^{n(\theta-1)+n}|B|^{\theta-1}\int_B\mathbf1_E(z)\,dz
 =r^\beta |B|^{\theta-1}|E\cap B|.
\]
这给出第一式；又 $|rE|^\theta=r^\beta|E|^\theta$，从而给出光晕恒等式。一般幂型阈值在伸缩下乘以 $r^{n\sigma}$，故与算子的齐次度相同当且仅当 $n\sigma=\beta$。
\end{proof}

\begin{proposition}[固定阈值的退化]
\label{fnd:raw-divergence}
对任意 $\lambda>0$ 以及
$\mathcal B\in\{\mathcal Q,\mathcal R,\mathcal Q_c\}$，
\[
 \sup_{0<|E|<\infty}
 \frac{|\{M_{\beta,\mathcal B}\mathbf1_E>\lambda\}|}{|E|}=\infty.
\]
\end{proposition}

\begin{proof}
令 $E_0=[-1/2,1/2]^n$。若 $\|x\|_\infty\leq R$，则以 $x$ 为中心、边长 $2R+1$ 的方体包含 $E_0$，所以对三种基均有
\[
 M_{\beta,\mathcal B}\mathbf1_{E_0}(x)\geq(2R+1)^{-(n-\beta)}.
\]
因此当 $u>0$ 充分小时，存在仅依赖于 $n,\beta$ 的 $c>0$ 使
\[
 |\{M_{\beta,\mathcal B}\mathbf1_{E_0}>u\}|
 \geq c u^{-n/(n-\beta)}.
\]
在伸缩式中取 $E=rE_0$，并令 $r\to\infty$，即得相应商至少为常数倍的
$\lambda^{-n/(n-\beta)}r^{\beta n/(n-\beta)}$，从而发散。
\end{proof}

\begin{theorem}[双密度刚性与最佳代数误差]
\label{fnd:double-density}
设 $0<|E|<\infty$，且 $B$ 为严格见证，即
\[
 |B|^{\theta-1}|E\cap B|>\alpha|E|^\theta.
\]
写
\[
 s:=|E\cap B|,\qquad a:=\frac{s}{|B|},\qquad b:=\frac{s}{|E|}.
\]
则
\[
 a^{1-\theta}b^\theta>\alpha,\qquad
 a>\alpha^{1/(1-\theta)},\qquad b>\alpha^{1/\theta},
\]
并且
\[
 \alpha^{1/\theta}|E|<|B|<\alpha^{-1/(1-\theta)}|E|.
\]
令
\[
 \rho_\theta(\alpha):=
 \max\left\{1-\alpha^{1/\theta},\,
 \alpha^{-1/(1-\theta)}-1\right\}.
\]
则
\[
 |B\triangle E|<\rho_\theta(\alpha)|E|.
\]
而且在只使用严格见证不等式及 $s\leq |B|,|E|$ 的意义下，该常数最优。
\end{theorem}

\begin{proof}
将严格见证不等式除以 $|B|^{1-\theta}|E|^\theta$，可得
\[
 a^{1-\theta}b^\theta>\alpha.
\]
因 $a,b\leq1$，分别去掉不超过 $1$ 的另一个因子即可得到两个密度下界。与 $s\leq|B|$、$s\leq|E|$ 结合，得到尺度窗。

为证明对称差估计，利用齐次性令 $|E|=1$，记 $r=|B|$。此时
$s>\alpha r^{1-\theta}$，且 $s\leq\min\{1,r\}$，故
\[
 \alpha^{1/\theta}<r<\alpha^{-1/(1-\theta)}.
\]
又
\[
 |B\triangle E|=1+r-2s
 <f(r):=1+r-2\alpha r^{1-\theta}.
\]
而 $f''(r)=2\alpha\theta(1-\theta)r^{-1-\theta}>0$。凸函数在闭区间的最大值由端点控制，故
\[
 |B\triangle E|<
 \max\{f(\alpha^{1/\theta}),f(\alpha^{-1/(1-\theta)})\}
 =\rho_\theta(\alpha).
\]
最后，令 $|E|=1$。取 $B\subset E$ 且 $|B|\downarrow\alpha^{1/\theta}$，或取 $E\subset B$ 且 $|B|\uparrow\alpha^{-1/(1-\theta)}$，均可保持严格见证并分别逼近上述两个端点误差。这说明统一常数不能减小。
\end{proof}

\begin{corollary}[端点展开]
\label{fnd:rho-expansion}
当 $t\downarrow0$ 时，
\[
 \rho_\theta(1-t)=
 \max\left\{\frac1\theta,\frac1{1-\theta}\right\}t+O_\theta(t^2).
\]
\end{corollary}

\begin{proof}
分别对 $(1-t)^{1/\theta}$ 与 $(1-t)^{-1/(1-\theta)}$ 作 Taylor 展开。
\end{proof}

\begin{remark}
\label{fnd:outer-weight}
后文若使用外层权常数，约定
\[
 C^w_{\beta,\mathcal B}(\alpha):=
 \sup_{0<|E|<\infty,\,0<w(E)<\infty}
 \frac{w(H_{\beta,\mathcal B}(\alpha,E))}{w(E)}.
\]
这里算子、阈值和集合尺度始终由 Lebesgue 测度决定；它不同于以内禀测度 $w(x)\,dx$ 定义的极大算子，也不同于一般的两权范数不等式。
\end{remark}

\section{近重合几何与端点阶}
\label{sec:geometry}

\begin{lemma}[近重合矩形]
\label{geo:rectangles}
设轴平行矩形 $R=\prod_{j=1}^nI_j$ 与 $R_0=\prod_{j=1}^nI_j^0$ 满足
\[
 |R\triangle R_0|\leq\varepsilon|R_0|,\qquad0<\varepsilon<1.
\]
对区间 $I$，用 $I[h]$ 表示把两端各延长 $h|I|$ 所得区间；令
\[
 A_n(\varepsilon):=\frac{1+\varepsilon}{(1-\varepsilon)^{n-1}},
 \qquad h_n(\varepsilon):=A_n(\varepsilon)-(1-\varepsilon).
\]
则 $R\subset R_0[h_n(\varepsilon)]$。特别地，对 $0<\varepsilon\leq1/2$，存在只依赖于 $n$ 的 $c_n$ 使 $h_n(\varepsilon)\leq c_n\varepsilon$。
\end{lemma}

\begin{proof}
写 $L_j=|I_j^0|$、$\ell_j=|I_j|$、$o_j=|I_j\cap I_j^0|$。由假设，
\[
 (1-\varepsilon)|R_0|\leq|R|\leq(1+\varepsilon)|R_0|,
 \qquad \prod_{j=1}^no_j=|R\cap R_0|\geq(1-\varepsilon)|R_0|.
\]
逐坐标地 $o_j\geq(1-\varepsilon)L_j$；再由 $\ell_i\geq o_i$ 得
\[
 \ell_j\leq\frac{(1+\varepsilon)\prod_iL_i}
 {\prod_{i\ne j}(1-\varepsilon)L_i}=A_n(\varepsilon)L_j.
\]
故 $I_j$ 在 $I_j^0$ 两端之外的长度均不超过
$\ell_j-o_j\leq h_n(\varepsilon)L_j$。逐坐标相乘即得所求包含。最后的线性界由 $h_n$ 在 $[0,1/2]$ 上的有界导数给出。
\end{proof}

\begin{lemma}[近重合中心方体]
\label{geo:centered-cubes}
设 $Q,Q_0$ 是边长分别为 $\ell,\ell_0$、中心分别为 $x,x_0$ 的轴平行方体，且
\[
 |Q\triangle Q_0|\leq\varepsilon|Q_0|,\qquad0<\varepsilon<\frac12.
\]
则
\[
 \|x-x_0\|_\infty\leq\frac32\varepsilon\ell_0.
\]
\end{lemma}

\begin{proof}
体积差控制给出 $(1-\varepsilon)\ell_0^n\leq\ell^n\leq(1+\varepsilon)\ell_0^n$。利用 Bernoulli 不等式的较弱版本可得
$(1-\varepsilon)\ell_0\leq\ell\leq(1+\varepsilon)\ell_0$。并且每个坐标投影的交长至少为 $(1-\varepsilon)\ell_0$，否则整个交集体积小于 $(1-\varepsilon)\ell_0^n$。

固定一个坐标。若两个投影区间相互包含，中心距离至多 $|\ell-\ell_0|/2\leq\varepsilon\ell_0/2$。否则其交长为 $(\ell+\ell_0)/2-|x_j-x_{0,j}|$，从而
\[
 |x_j-x_{0,j}|\leq\frac{\ell+\ell_0}{2}-(1-\varepsilon)\ell_0
 \leq\frac32\varepsilon\ell_0.
\]
对所有坐标取最大值即可。
\end{proof}

\begin{theorem}[三类基的显式端点界]
\label{geo:explicit-bounds}
令
\[
 L_\alpha:=\alpha^{-1/(n-\beta)},\qquad
 A_\alpha:=\alpha^{-n/(n-\beta)},\qquad
 B_\alpha:=\alpha^{1/\beta},
\]
并令 $\rho=\rho_\theta(\alpha)$、
$\varepsilon_\alpha:=2\rho\alpha^{-1/\theta}$。当
$\varepsilon_\alpha<1/2$ 时，
\begin{equation}
 C_{\beta,\mathcal Q}(\alpha),C_{\beta,\mathcal R}(\alpha)
 \leq (1+2h_n(\varepsilon_\alpha))^n
 \alpha^{-1/(1-\theta)},
 \label{geo:noncentered-upper}
\end{equation}
及
\begin{equation}
 C_{\beta,\mathcal Q_c}(\alpha)
 \leq(3\varepsilon_\alpha)^n\alpha^{-1/(1-\theta)}.
 \label{geo:centered-upper}
\end{equation}
对每个 $0<\alpha<1$，还有
\[
 C_{\beta,\mathcal Q}(\alpha)\geq(2L_\alpha-1)^n,
 \quad C_{\beta,\mathcal R}(\alpha)\geq1+2n(A_\alpha-1),
\]
以及
\[
 C_{\beta,\mathcal Q_c}(\alpha)
 \geq\max\{(1-B_\alpha)^n,(L_\alpha-1)^n\}.
\]
\end{theorem}

\begin{proof}
固定 $E$，记 $H=H_{\beta,\mathcal B}(\alpha,E)$。若 $H$ 非空，为每个 $x\in H$ 选择严格见证 $B_x$，再固定一个 $B_0$。由定理~\ref{fnd:double-density}，
\[
 |B_x\triangle B_0|<2\rho|E|
 <\varepsilon_\alpha|B_0|,
\]
其中最后一步用了 $|B_0|>\alpha^{1/\theta}|E|$。当
$\mathcal B=\mathcal Q$ 或 $\mathcal R$ 时，引理~\ref{geo:rectangles} 给出
$H\subset B_0[h_n(\varepsilon_\alpha)]$，于是
\[
 \frac{|H|}{|E|}
 \leq(1+2h_n(\varepsilon_\alpha))^n\frac{|B_0|}{|E|}
 \leq(1+2h_n(\varepsilon_\alpha))^n\alpha^{-1/(1-\theta)}.
\]
当 $\mathcal B=\mathcal Q_c$ 时，$x$ 是 $B_x$ 的中心。由引理~\ref{geo:centered-cubes}，所有这些中心落在以 $B_0$ 中心为中心、边长 $3\varepsilon_\alpha\ell(B_0)$ 的方体中，给出第二个上界。

下界均取方体测试集。对 $\mathcal Q$，令 $E=[0,1]^n$。若 $x\in(1-L_\alpha,L_\alpha)^n$，则存在一个包含 $E$ 与 $x$、边长严格小于 $L_\alpha$ 的方体，故该点属于光晕。对 $\mathcal R$，同样取 $E=[0,1]^n$；沿任一坐标轴在两侧附加宽度 $A_\alpha-1$ 的开薄层，用包含 $E$ 和该点的矩形作严格见证。这些薄层模零不交，给出第二个下界。

最后令 $E=[-1/2,1/2]^n$。以 $x$ 为中心、完全包含在 $E$ 内的方体给出边长 $1-B_\alpha$ 的中心方体；以 $x$ 为中心、完全捕获 $E$ 的最小方体给出边长 $L_\alpha-1$ 的中心方体。两个严格包含关系分别给出两个候选集；它们同心，故其并的体积是所列最大值。
\end{proof}

\begin{corollary}[端点幂指数与可验证的首项下界]
\label{geo:endpoint-orders}
当 $\alpha\uparrow1$ 时，
\[
 C_{\beta,\mathcal Q}(\alpha)-1\asymp_{n,\beta}1-\alpha,
 \qquad C_{\beta,\mathcal R}(\alpha)-1\asymp_{n,\beta}1-\alpha,
\]
以及
\[
 C_{\beta,\mathcal Q_c}(\alpha)\asymp_{n,\beta}(1-\alpha)^n.
\]
更精确地，
\[
 \liminf_{\alpha\uparrow1}
 \frac{C_{\beta,\mathcal Q}(\alpha)-1}{1-\alpha}\geq\frac{2n}{n-\beta},
 \quad
 \liminf_{\alpha\uparrow1}
 \frac{C_{\beta,\mathcal R}(\alpha)-1}{1-\alpha}\geq\frac{2n^2}{n-\beta},
\]
并且
\[
 \liminf_{\alpha\uparrow1}
 \frac{C_{\beta,\mathcal Q_c}(\alpha)}{(1-\alpha)^n}
 \geq\max\{\beta^{-n},(n-\beta)^{-n}\}.
\]
\end{corollary}

\begin{proof}
写 $t=1-\alpha$。有
\[
 L_\alpha-1=\frac{t}{n-\beta}+O_{n,\beta}(t^2),\quad
 A_\alpha-1=\frac{n}{n-\beta}t+O_{n,\beta}(t^2),\quad
 1-B_\alpha=\frac{t}{\beta}+O_\beta(t^2).
\]
代入定理~\ref{geo:explicit-bounds} 的下界给出首项下界。上界来自推论~\ref{fnd:rho-expansion}、$h_n(u)=O_n(u)$ 以及定理~\ref{geo:explicit-bounds}。
\end{proof}

\begin{remark}
\label{geo:scope}
本节的 ``sharp'' 仅指幂指数不能改善。中心高维的全体集合首项常数不由上述端点测试或近重合估计决定；后文的固定方体切向极限只能先给出相应的全体常数下界。
\end{remark}

\section{中心方体光晕的精确切向多面体}
\label{tan:section}

本节固定 $0<\beta<n$，并令
\[
 Q_0=[-1/2,1/2]^n,\qquad c=\min\{\beta,n-\beta\},
 \qquad \sigma=2c,\qquad 0<\sigma\le n.
\]
对中心轴平行方体基 $\mathcal Q_c$，记
\[
 H_{\beta,\mathcal Q_c}(\alpha,E)
 =\left\{x:\sup_{\ell>0}\ell^{\beta-n}
   |E\cap (x+[-\ell/2,\ell/2]^n)|>\alpha |E|^{\beta/n}\right\}.
\]
这里的严格不等式与下文所有边界论证一致。

令 $k=\lfloor\sigma\rfloor$；若 $\sigma=n$，仍取 $k=n$。给定
$v\in\mathbb R^n$，将 $|v_1|,\ldots,|v_n|$ 按递减顺序记为
$v_1^*\ge\cdots\ge v_n^*$，并定义
\begin{equation}
 N_c(v)=
 \begin{cases}
 \displaystyle\sum_{j=1}^{k}v_j^*+(\sigma-k)v_{k+1}^*,&k<n,\\[4pt]
 \displaystyle\sum_{j=1}^{n}v_j^*,&k=n.
 \end{cases}
 \label{tan:normal-form}
\end{equation}
当 $k=0$ 时，空和为零，所以 $N_c(v)=\sigma v_1^*$。令
\begin{equation}
 S_{n,c}=\{v\in\mathbb R^n:N_c(v)<1\}.
 \label{tan:polytope}
\end{equation}

\begin{lemma}[缩放参数的精确消去]
\label{tan:parameter-elimination}
对每个 $v\in\mathbb R^n$，有
\begin{equation}
 N_c(v)=\inf_{u\ge0}\left\{cu+
 \sum_{j=1}^{n}\bigl(|v_j|-u/2\bigr)_+\right\}.
 \label{tan:infimum}
\end{equation}
\end{lemma}

\begin{proof}
令 $z=u/2$。右端变为
\[
 G(z)=\sigma z+\sum_{j=1}^{n}(v_j^*-z)_+\qquad(z\ge0).
\]
在不含任何 $v_j^*$ 的开区间内，若恰有 $r$ 个坐标大于 $z$，则
$G'(z)=\sigma-r$。若 $k<\sigma<k+1$，最小值在
$z=v_{k+1}^*$ 取得；若 $\sigma=k<n$，最小区间为
$[v_{k+1}^*,v_k^*]$。代入即得第一种情形。若 $\sigma=n$，
$G$ 在 $[0,v_n^*]$ 上恒等于 $\sum_jv_j^*$。这也给出第二种情形。
\end{proof}

\begin{theorem}[单位方体的精确切向极限]
\label{tan:unit-cube-limit}
当 $t\downarrow0$ 时，
\begin{equation}
 t^{-n}\bigl|H_{\beta,\mathcal Q_c}(1-t,Q_0)\bigr|
 \longrightarrow |S_{n,c}|
 =\frac{2^n}{k!\,(2c)^{n-k}}.
 \label{tan:cube-limit}
\end{equation}
\end{theorem}

\begin{proof}
先固定 $v\in\mathbb R^n$ 与 $u\ge0$。对外扩方体
\[
 Q_t^+=tv+[-(1+tu)/2,(1+tu)/2]^n
\]
逐坐标计算交长，并在 $u$、$\|v\|_\infty$ 有界时作一致展开，得到
\begin{equation}
 \log\!\left((1+tu)^{\beta-n}|Q_0\cap Q_t^+|\right)
 =-t\left((n-\beta)u+
 \sum_{j=1}^{n}(|v_j|-u/2)_+\right)+O(t^2).
 \label{tan:outer-expansion}
\end{equation}
对内缩方体
\[
 Q_t^-=tv+[-(1-tu)/2,(1-tu)/2]^n
\]
同样有
\begin{equation}
 \log\!\left((1-tu)^{\beta-n}|Q_0\cap Q_t^-|\right)
 =-t\left(\beta u+
 \sum_{j=1}^{n}(|v_j|-u/2)_+\right)+O(t^2).
 \label{tan:inner-expansion}
\end{equation}
两式的余项在任意固定紧参数集上一致。由于
$\log(1-t)=-t+O(t^2)$，若右端括号中的量严格小于 $1$，相应方体便是
$tv\in H_{\beta,\mathcal Q_c}(1-t,Q_0)$ 的严格见证。由
引理~\ref{tan:parameter-elimination}，这说明对每个 $v\in S_{n,c}$，
充分小的 $t$ 都满足
\begin{equation}
 tv\in H_{\beta,\mathcal Q_c}(1-t,Q_0).
 \label{tan:inner-inclusion}
\end{equation}

反过来，设 $t_r\downarrow0$，$x_r\in
H_{\beta,\mathcal Q_c}(1-t_r,Q_0)$，并以边长 $\ell_r$ 的中心方体为严格见证。
记 $s_r=|Q_0\cap Q_r|$。见证不等式以及 $s_r\le\min\{1,\ell_r^n\}$ 先给出
\[
 (1-t_r)^{1/\beta}<\ell_r<(1-t_r)^{-1/(n-\beta)},
 \qquad s_r=1+O(t_r).
\]
故 $|Q_r\triangle Q_0|=O(t_r)$。对两个轴平行方体逐坐标比较交长可知
$\ell_r=1+O(t_r)$ 与 $|x_r|=O(t_r)$。特别地，存在只依赖于 $n,\beta$ 的
$R$，使
\begin{equation}
 t^{-1}H_{\beta,\mathcal Q_c}(1-t,Q_0)\subset[-R,R]^n
 \quad\text{当 $t$ 充分小时}.
 \label{tan:uniform-support}
\end{equation}

若 $v_r=x_r/t_r\to v$，再取子列可设
$a_r=(\ell_r-1)/t_r\to a$ 且 $a_r$ 同号。若 $a_r\ge0$，
把式~\eqref{tan:outer-expansion}用于 $u=a_r$；若 $a_r\le0$，把
式~\eqref{tan:inner-expansion}用于 $u=-a_r$。由严格见证不等式取极限，分别得到
\[
 (n-\beta)a+\sum_j(|v_j|-a/2)_+\le1
 \quad(a\ge0),
\]
或
\[
 \beta(-a)+\sum_j(|v_j|-(-a)/2)_+\le1
 \quad(a\le0).
\]
由引理~\ref{tan:parameter-elimination}，两种情形都蕴含 $N_c(v)\le1$。
因此，若 $N_c(v)>1$，则 $v\notin t^{-1}H_{\beta,\mathcal Q_c}(1-t,Q_0)$
对充分小的 $t$ 成立。结合式~\eqref{tan:inner-inclusion}，缩放光晕的示性函数在
$\partial S_{n,c}$ 外逐点收敛到 $\mathbf1_{S_{n,c}}$。集合 $S_{n,c}$ 是有界多面体的内部，
其边界为零测集；再用式~\eqref{tan:uniform-support}与支配收敛定理，得到
\[
 t^{-n}|H_{\beta,\mathcal Q_c}(1-t,Q_0)|\to |S_{n,c}|.
\]

最后计算体积。先在 $v_1\ge\cdots\ge v_n\ge0$ 的锥内积分。若 $k<n$，置
$z=v_{k+1}$、$y_j=v_j-z$（$1\le j\le k$）。尾部变量的体积为
$z^{n-k-1}/(n-k-1)!$，而有序单纯形
$y_1\ge\cdots\ge y_k\ge0$、$\sum_jy_j<A$ 的体积为 $A^k/(k!)^2$。
于是该锥内的体积为
\[
 \frac1{(n-k-1)!(k!)^2}
 \int_0^{1/\sigma}z^{n-k-1}(1-\sigma z)^k\,dz
 =\frac1{k!\,n!\,\sigma^{n-k}}.
\]
乘以 $2^n n!$ 即得所示结果。$k=n$ 时 $S_{n,c}$ 是 $\ell^1$ 单位球，体积
$2^n/n!$，与同一公式一致。
\end{proof}

\begin{corollary}[全体中心方体常数的严格下界]
\label{tan:global-liminf}
若
\[
 C_{\beta,\mathcal Q_c}(\alpha)=
 \sup_{0<|E|<\infty}\frac{|H_{\beta,\mathcal Q_c}(\alpha,E)|}{|E|},
\]
则
\begin{equation}
 \liminf_{\alpha\uparrow1}
 \frac{C_{\beta,\mathcal Q_c}(\alpha)}{(1-\alpha)^n}
 \ge\frac{2^n}{k!\,[2\min\{\beta,n-\beta\}]^{n-k}}.
 \label{tan:global-liminf-formula}
\end{equation}
\end{corollary}

\begin{proof}
在常数的上确界中取 $E=Q_0$，并应用定理~\ref{tan:unit-cube-limit}。
\end{proof}

\begin{remark}[首项的准确状态]
\label{tan:status}
推论~\ref{tan:global-liminf}只是由一个固定测试集合给出的下界。它既不证明极限存在，
也不证明该下界等于全体可测集的最优首项。尤其当 $c>1/2$ 时，右端严格大于旧候选
$c^{-n}$，因而该旧候选不可能是正确首项。尚待解决的问题是：全体集合的最优一阶
形状变分是否恰由 $Q_0$ 实现。
\end{remark}

\section{固定凸体与卷积体}
\label{cvx:section}

\subsection{中心对称凸体的端点阶}

设 $K\subset\mathbb R^n$ 是有界中心对称凸体，且 $0\in\operatorname{int}K$。定义
\[
 \mathcal B^K_{\mathrm u}(x)=\{z+rK:r>0,\ x\in z+rK\},
 \qquad
 \mathcal B^K_{\mathrm c}(x)=\{x+rK:r>0\}.
\]
相应的分数次算子与光晕沿用正文的尺度不变阈值
$\alpha|E|^{\beta/n}$。以下近重合引理是把方体几何替换为固定凸体几何所需的唯一
定量输入。

\begin{lemma}[近重合位似凸体]
\label{cvx:near-homothetic}
设 $K\subset\mathbb R^n$ 是有界满维凸体，且 $0\in\operatorname{int}K$。存在只依赖于
$K$ 的常数 $\varepsilon_K,C_K>0$，使得若
\[
 B=x+rK,\qquad B_0=x_0+r_0K,\qquad
 |B\triangle B_0|\le\varepsilon|B_0|,\qquad 0<\varepsilon<\varepsilon_K,
\]
则
\begin{equation}
 \left|\frac r{r_0}-1\right|\le C_K\varepsilon,
 \qquad
 |x-x_0|_2\le C_K\varepsilon r_0.
 \label{cvx:near-homothetic-bound}
\end{equation}
\end{lemma}

\begin{proof}
记 $V=|K|$。由体积差不超过对称差，
\[
 |r^n-r_0^n|V\le\varepsilon r_0^nV.
\]
取 $\varepsilon_K\le1/2$。函数 $s\mapsto s^{1/n}$ 在包含
$[(1-\varepsilon_K),(1+\varepsilon_K)]$ 的紧区间上 Lipschitz，故第一个结论成立。
再由三角不等式，
\begin{equation}
 |(x+r_0K)\triangle(x_0+r_0K)|
 \le |B\triangle B_0|+|(x+rK)\triangle(x+r_0K)|
 \le2\varepsilon r_0^nV.
 \label{cvx:translation-reduction}
\end{equation}

为控制平移，取 $0<r_*<R_*$ 使
$r_*B_2^n\subset K\subset R_*B_2^n$。对每个单位向量 $e$，沿
$\mathbb Re$ 作切片。若 $y\in(r_*/2)B_2^{n-1}\subset e^\perp$，
则 $K$ 在 $y+\mathbb Re$ 上的截段长度至少 $\sqrt3r_*$。所以当
$0<s<\sqrt3r_*$ 时，Fubini 定理给出
\begin{equation}
 |(K+se)\triangle K|\ge2s\,\omega_{n-1}(r_*/2)^{n-1}.
 \label{cvx:translation-lower}
\end{equation}
另一方面，在环带 $\{h:\sqrt3r_*/2\le|h|_2\le2R_*\}$ 上，连续函数
$h\mapsto|(K+h)\triangle K|$ 有严格正的最小值；而 $|h|_2>2R_*$ 时两个体不交。
故当右端足够小时，式~\eqref{cvx:translation-reduction} 迫使
$h=(x-x_0)/r_0$ 落入式~\eqref{cvx:translation-lower} 的适用范围，进而
$|h|_2\le C\varepsilon$，这就得到第二个结论。该证明未使用 $K=-K$。
\end{proof}

\begin{theorem}[固定中心对称凸体的端点阶]
\label{cvx:endpoint-order}
对每个固定的 $K$ 与 $0<\beta<n$，当 $\alpha\uparrow1$ 时，
\begin{equation}
 C_{\beta,\mathcal B^K_{\mathrm u}}(\alpha)-1\asymp_{K,\beta}1-\alpha,
 \qquad
 C_{\beta,\mathcal B^K_{\mathrm c}}(\alpha)
 \asymp_{K,\beta}(1-\alpha)^n.
 \label{cvx:endpoint-order-formula}
\end{equation}
更具体地，令
\[
 B_\alpha=\alpha^{1/\beta},\qquad
 L_\alpha=\alpha^{-1/(n-\beta)},\qquad\theta=\beta/n,
\]
则
\begin{equation}
 C_{\beta,\mathcal B^K_{\mathrm u}}(\alpha)\ge(2L_\alpha-1)^n,
 \qquad
 C_{\beta,\mathcal B^K_{\mathrm c}}(\alpha)
 \ge\max\{(1-B_\alpha)^n,(L_\alpha-1)^n\}.
 \label{cvx:endpoint-lower}
\end{equation}
\end{theorem}

\begin{proof}
只证上界，因为下界由 $E=K$ 直接给出。固定一个严格见证
$B_0=x_0+r_0K$。双密度估计给出：当 $\alpha$ 接近 $1$ 时，任何严格见证
$B_x=z_x+r_xK$ 满足
\[
 |B_x\triangle B_0|\le C_{\beta}(1-\alpha)|B_0|.
\]
对引理~\ref{cvx:near-homothetic} 应用此式，可得
\[
 |r_x/r_0-1|+\|z_x-x_0\|_K/r_0\le C_{K,\beta}(1-\alpha).
\]
非中心情形中 $x\in B_x$，故光晕被一个半径为
$(1+C_{K,\beta}(1-\alpha))r_0$ 的固定位似体包含；再用见证的尺度窗比较
$|B_0|$ 与 $|E|$，得到线性上界。中心情形中 $x=z_x$，故光晕包含在
$x_0+C_{K,\beta}(1-\alpha)r_0K$ 内，从而得到 $n$ 次上界。

为证式~\eqref{cvx:endpoint-lower}，取 $E=K$。若 $1<r<L_\alpha$ 且
$z\in(r-1)K$，则 $K\subset z+rK$；这些见证的并包含 $(2r-1)K$。
令 $r\uparrow L_\alpha$ 得到非中心下界。中心情形中，若 $B_\alpha<r<1$ 且
$x\in(1-r)K$，则 $x+rK\subset K$；若 $1<r<L_\alpha$ 且
$x\in(r-1)K$，则 $K\subset x+rK$。令参数趋近端点，即得所示两个下界。
\end{proof}

\subsection{卷积体与投影体的固定尺度下界}

以下部分不要求 $K$ 中心对称。选择参考点使 $0\in\operatorname{int}K$，定义锚定的
同似基
\[
 M^K_{\beta,\mathrm c}f(x)=\sup_{s>0}|sK|^{\beta/n-1}
 \int_{x+sK}|f(y)|\,dy,
\]
并记相应光晕为 $H^K_{\beta,\mathrm c}(\alpha,E)$。协方差函数与卷积体为
\[
 g_K(x)=|K\cap(K+x)|,
 \qquad C^K_\alpha=\{x:g_K(x)>\alpha|K|\}.
\]

\begin{proposition}[卷积体嵌入与精确同似表示]
\label{cvx:convolution-embedding}
对 $0<\alpha<1$，有
\begin{equation}
 C^K_\alpha\subset H^K_{\beta,\mathrm c}(\alpha,K),
 \label{cvx:embedding}
\end{equation}
并且
\begin{equation}
 H^K_{\beta,\mathrm c}(\alpha,K)
 =\bigcup_{s>0}\left\{x:|K\cap(x+sK)|>\alpha s^{n-\beta}|K|\right\}.
 \label{cvx:homothetic-representation}
\end{equation}
\end{proposition}

\begin{proof}
将 $E=K$ 代入定义，并除以正数 $|K|^{\beta/n-1}$，即得
\[
 |K\cap(x+sK)|>\alpha s^{n-\beta}|K|.
\]
对 $s$ 取并便得到式~\eqref{cvx:homothetic-representation}；取 $s=1$ 即得
式~\eqref{cvx:embedding}。
\end{proof}

令 $\Pi K$ 为投影体，其支持函数规范化为
\[
 h_{\Pi K}(v)=|v|_2\,\bigl|P_{v^\perp}K\bigr|_{n-1}\quad(v\ne0),
 \qquad h_{\Pi K}(0)=0,
\]
并写 $\Pi^*K=(\Pi K)^\circ$。

\begin{lemma}[卷积体的近端极限]
\label{cvx:covariogram-limit}
当 $\alpha\uparrow1$ 时，
\begin{equation}
 (1-\alpha)^{-1}C^K_\alpha\longrightarrow |K|\Pi^*K
 \quad\text{在 Hausdorff 距离下},
 \label{cvx:hausdorff-limit}
\end{equation}
因而
\begin{equation}
 \lim_{\alpha\uparrow1}\frac{|C^K_\alpha|}{(1-\alpha)^n}
 =|K|^n|\Pi^*K|.
 \label{cvx:covariogram-volume}
\end{equation}
\end{lemma}

\begin{proof}
固定单位向量 $u$，并沿平行于 $u$ 的直线切片。若 $\ell_u(y)$ 是
$K$ 在 $y+\mathbb Ru$ 上的截段长度，则
\[
 g_K(tu)=\int_{P_{u^\perp}K}(\ell_u(y)-t)_+\,dy.
\]
所以由支配收敛定理，
\begin{equation}
 \frac{|K|-g_K(tu)}t\longrightarrow |P_{u^\perp}K|_{n-1}=h_{\Pi K}(u).
 \label{cvx:directional-derivative}
\end{equation}
齐次性给出 $g_K(tv)=|K|-t h_{\Pi K}(v)+o(t)$。

为把逐方向结论提升到 Hausdorff 收敛，取 $r_*>0$ 使某个平移的 $r_*B_2^n$ 包含于
$K$。前述切片公式在单位球面上一致给出常数 $0<m\le M<\infty$，使小 $t>0$ 时
\[
 mt\le |K|-g_K(tu)\le Mt\qquad(u\in S^{n-1}).
\]
故 $(1-\alpha)^{-1}C^K_\alpha$ 被两个固定欧氏球夹住。又由
$g_K^{1/n}$ 的凹性，$C^K_\alpha$ 的闭包为凸体；其径向函数因而在球面上一致
Lipschitz。式~\eqref{cvx:directional-derivative} 给出径向函数逐点收敛到
$|K|/h_{\Pi K}(u)$。有限网逼近将其提升为一致收敛，得到
式~\eqref{cvx:hausdorff-limit}。凸体体积对 Hausdorff 收敛连续，故
式~\eqref{cvx:covariogram-volume} 成立。
\end{proof}

\begin{corollary}[投影体给出的中心光晕下界]
\label{cvx:projection-liminf}
有
\begin{equation}
 \liminf_{\alpha\uparrow1}
 \frac{|H^K_{\beta,\mathrm c}(\alpha,K)|}{|K|(1-\alpha)^n}
 \ge |K|^{n-1}|\Pi^*K|.
 \label{cvx:projection-liminf-formula}
\end{equation}
\end{corollary}

\begin{proof}
由式~\eqref{cvx:embedding}，光晕测度至少为 $|C^K_\alpha|$。除以
$|K|(1-\alpha)^n$ 后使用引理~\ref{cvx:covariogram-limit}。
\end{proof}

\section{一维精确常数}\label{sec:one-dimensional}

本节固定 $n=1$、$0<\beta<1$，并记
\[
 q:=\frac{1}{1-\beta},\qquad
 A_\beta(t):=1-(1-t)^{1/\beta},\qquad
 D_\beta(t):=(1-t)^{-1/(1-\beta)}-1.
\]
这里 $t=1-\alpha$。为避免混淆，$A_\beta$ 对应完全包含于资料集的中心见证，
$D_\beta$ 对应完整捕获资料集的中心见证；它们并不是两个可同时相加的光晕部分。

\begin{lemma}[两两相交区间族]\label{one:pairwise}
设 $\mathscr J$ 是非空区间族，且任意 $I,J\in\mathscr J$ 相交。则
$U:=\bigcup_{I\in\mathscr J}I$ 是区间。若还存在 $L<\infty$ 使
\[
 |I\cup J|\le L\qquad(I,J\in\mathscr J),
\]
则 $|U|\le L$。
\end{lemma}

\begin{proof}
两两相交的区间并是区间。取 $I_0\in\mathscr J$。由
$|I\cup I_0|\le L$ 可知所有端点有界。令 $a:=\inf U$、$b:=\sup U$。
给定 $\varepsilon>0$，可取 $I_-,I_+\in\mathscr J$ 满足
\[
 \inf I_-<a+\varepsilon,\qquad \sup I_+>b-\varepsilon.
\]
二者相交，故 $I_-\cup I_+$ 为区间，从而
$b-a\le |I_-\cup I_+|+2\varepsilon\le L+2\varepsilon$。令
$\varepsilon\downarrow0$ 即得结论。
\end{proof}

\begin{theorem}[一维非中心精确常数]\label{one:uncentered}
若 $2^{-\beta}<\alpha<1$，则
\[
 C_{\beta,\mathcal Q}(\alpha)
 =C_{\beta,\mathcal R}(\alpha)
 =2\alpha^{-1/(1-\beta)}-1.
\]
\end{theorem}

\begin{proof}
在一维，$\mathcal Q=\mathcal R$ 是全体有界区间。由尺度不变性，只须处理
$|E|=1$。令 $I$ 为严格见证区间，并写 $s_I:=|E\cap I|$。见证不等式给出
\[
 |I|<\left(\frac{s_I}{\alpha}\right)^q,
 \qquad s_I>\alpha^{1/\beta}>\frac12,
 \qquad \alpha^{1/\beta}<s_I\le1.
\]
后两个不等式分别由 $s_I\le |I|$、$s_I\le |E|$ 与严格见证条件得到。
因此两个严格见证区间 $I,J$ 满足
\[
 |E\cap I\cap J|\ge s_I+s_J-|E|=s_I+s_J-1>0.
\]
令 $g(s):=(s/\alpha)^q-s$。在 $[\alpha^{1/\beta},1]$ 上，
\[
 g'(s)=q\alpha^{-q}s^{q-1}-1\ge q-1>0,
\]
故 $g(s)\le g(1)=\alpha^{-q}-1$。又因 $I\cap J$ 的测度至少
$s_I+s_J-1$，有
\[
 |I\cup J|
 <\left(\frac{s_I}{\alpha}\right)^q+
   \left(\frac{s_J}{\alpha}\right)^q-(s_I+s_J-1)
 =1+g(s_I)+g(s_J)
 \le2\alpha^{-q}-1.
\]
引理~\ref{one:pairwise} 遂给出光晕的上界。

反向取 $E=[0,1]$。若 $x\in(1-\alpha^{-q},\alpha^{-q})$，则 $E$ 和 $x$
包含于长度严格小于 $\alpha^{-q}$ 的区间。该区间完整捕获 $E$，故严格见证
$x\in H_{\beta,\mathcal Q}(\alpha,E)$。该开区间的长度是
$2\alpha^{-q}-1$，给出相反方向的不等式。
\end{proof}

\begin{theorem}[一维中心精确常数]\label{one:centered}
若 $2^{-\beta}<\alpha<1$，则
\[
 C_{\beta,\mathcal Q_c}(\alpha)
 =\max\left\{1-\alpha^{1/\beta},\
 \alpha^{-1/(1-\beta)}-1\right\}.
\]
\end{theorem}

\begin{proof}
仍令 $|E|=1$。若中心严格见证区间 $I$ 的中心、长度和捕获量分别为
$c_I,\ell_I,s_I$，则 $s_I\in(B,1]$，其中
\[
 B:=\alpha^{1/\beta}>\frac12,\qquad L:=\alpha^{-q}.
\]
上个证明已经说明任意两个严格见证区间相交。对于这样的 $I,J$，有
\[
 |c_I-c_J|
 \le\frac{\ell_I+\ell_J}{2}-|I\cap J|
 <1+h(s_I)+h(s_J),
 \qquad h(s):=\frac12\left(\frac{s}{\alpha}\right)^q-s.
\]
函数 $h$ 在 $[B,1]$ 上凸，故其最大值在端点取得；而
\[
 h(B)=-\frac{B}{2},\qquad h(1)=\frac L2-1.
\]
于是任意两个光晕点的距离严格小于
\[
 1+2\max\left\{-\frac B2,\frac L2-1\right\}
 =\max\{1-B,L-1\}.
\]
实线上集合的测度不超过其直径，故得到上界。

再次取 $E=[0,1]$。若 $x\in(B/2,1-B/2)$，可选择以 $x$ 为中心、完全包含于
$E$ 且长度严格大于 $B$ 的区间，给出长度 $1-B$ 的内部光晕。若
$|x-1/2|<(L-1)/2$，可选择以 $x$ 为中心、完整包含 $E$ 且长度严格小于 $L$ 的区间，
给出长度 $L-1$ 的外部光晕。这两个区间集同心，故其并的长度是
$\max\{1-B,L-1\}$，上界达到。
\end{proof}

\begin{remark}[一阶相变]\label{one:first-order-phase}
定理~\ref{one:centered} 表明，当 $\beta<1/2$ 时，内含机制控制一阶项；
当 $\beta>1/2$ 时，完整捕获机制控制一阶项。在 $\beta=1/2$，两者一阶项相同，
但外部机制在二阶已经严格较大。
\end{remark}

\section{外层权、集中模量与匹配奇性模型}\label{sec:weights}

本节的权始终只作用于光晕集合的外层测度；极大算子及其阈值仍使用 Lebesgue 测度。设
\(0<\beta<n\)、\(\theta=\beta/n\)，并令 \(w\geq0\) 为局部可积且不恒为零的函数，
\[
 w(A):=\int_Aw(x)\,dx.
\]
对任一基 \(\mathcal B\)，记
\[
 C^w_{\beta,\mathcal B}(\alpha)
 :=\sup_{\substack{0<|E|<\infty\\0<w(E)<\infty}}
 \frac{w\bigl(H_{\beta,\mathcal B}(\alpha,E)\bigr)}{w(E)}.
\]
这里 \(H_{\beta,\mathcal B}(\alpha,E)\) 是前文定义的、以
\(\alpha|E|^\theta\) 为阈值的 Lebesgue 光晕。特别地，本节的结论不涉及把
\(w\,dx\) 放入平均式的内禀加权极大算子。

\begin{definition}\label{wgt:definition-concentration}
对 \(\mathcal D\in\{\mathcal Q,\mathcal R\}\) 及 \(0<u<1\)，定义外层集中模量
\[
 \Omega_{\mathcal D,w}(u)
 :=\sup_{\substack{D\in\mathcal D,\ w(D)>0\\ F\subset D,\ |F|\leq u|D|}}
 \frac{w(F)}{w(D)}.
\]
\end{definition}

为固定术语，对 \(1<p<\infty\) 定义
\[
 [w]_{A_p(\mathcal Q)}
 :=\sup_{Q\in\mathcal Q}
 \left(\frac1{|Q|}\int_Qw\right)
 \left(\frac1{|Q|}\int_Qw^{-1/(p-1)}\right)^{p-1},
\]
并令 \(A_\infty=\bigcup_{p>1}A_p(\mathcal Q)\)。将上式中的方体 \(Q\) 全部替换为轴平行矩形
\(R\in\mathcal R\)，所得的类记为 \(A_p^\ast\)，并令
\(A_\infty^\ast=\bigcup_{p>1}A_p^\ast\)。我们使用下列标准集中刻画：
\[
 w\in A_\infty\Longleftrightarrow\Omega_{\mathcal Q,w}(u)\longrightarrow0,
 \qquad
 w\in A_\infty^\ast\Longleftrightarrow\Omega_{\mathcal R,w}(u)\longrightarrow0,
 \quad(u\downarrow0),
\]
并且这些等价条件等同于相应基上存在 \(A<\infty\)、\(\delta>0\) 的幂型集中界。
这些标准事实及其加权 Solyanik 背景可参见
\cite{HagelsteinParissis2016,HagelsteinParissis2018}。

以下使用前文已证的见证聚拢结论。若 \(\rho_\theta(\alpha)\) 为双密度误差，且
\(\alpha\) 充分接近 \(1\)，则每个非空光晕可任选一个严格见证 \(B_0\) 作为锚点：
非中心方体或强矩形的全体见证均落入 \(B_0\) 的一个相对边界层中；中心方体的全体光晕
则落入与 \(B_0\) 同心、边长比为 \(O_{n,\beta}(1-\alpha)\) 的方体中。该结论同时给出
\begin{equation}
 \frac{|B_0\setminus E|}{|B_0|}\leq C_{n,\beta}\rho_\theta(\alpha),
 \qquad \rho_\theta(\alpha)\asymp_\theta1-\alpha.
\label{wgt:anchor-estimates}
\end{equation}

\begin{theorem}[集中模量控制]\label{wgt:theorem-forward}
存在仅依赖 \(n,\beta\) 的常数 \(a_{n,\beta}>0\)，使得当 \(\alpha\) 充分接近 \(1\) 时，
\begin{equation}
 C^w_{\beta,\mathcal Q}(\alpha)-1
 \leq4\Omega_{\mathcal Q,w}\bigl(a_{n,\beta}\rho_\theta(\alpha)\bigr),
\label{wgt:Q-forward}
\end{equation}
\begin{equation}
 C^w_{\beta,\mathcal R}(\alpha)-1
 \leq4\Omega_{\mathcal R,w}\bigl(a_{n,\beta}\rho_\theta(\alpha)\bigr),
\label{wgt:R-forward}
\end{equation}
并且
\begin{equation}
 C^w_{\beta,\mathcal Q_c}(\alpha)
 \leq
 \frac{\Omega_{\mathcal Q,w}\bigl(a_{n,\beta}\rho_\theta(\alpha)^n\bigr)}
 {1-\Omega_{\mathcal Q,w}\bigl(a_{n,\beta}\rho_\theta(\alpha)\bigr)}.
\label{wgt:Qc-forward}
\end{equation}
前两式在右端集中模量不超过 \(1/2\) 时成立；第三式要求其分母为正。
\end{theorem}

\begin{proof}
固定 \(E\)，并设 \(H=H_{\beta,\mathcal D}(\alpha,E)\) 非空。先取
\(\mathcal D\in\{\mathcal Q,\mathcal R\}\)，选定严格见证 \(B_0\)。由见证聚拢结论，存在
\(B_0^\ast\in\mathcal D\) 使 \(H\subset B_0^\ast\)，且在增大只依赖 \(n,\beta\) 的常数后，
\begin{equation}
 \frac{|B_0^\ast\setminus B_0|}{|B_0^\ast|}
 \leq a_{n,\beta}\rho_\theta(\alpha),\qquad
 \frac{|B_0^\ast\setminus E|}{|B_0^\ast|}
 \leq a_{n,\beta}\rho_\theta(\alpha),
\label{wgt:outer-errors}
\end{equation}
并由 \eqref{wgt:anchor-estimates} 有
\begin{equation}
 \frac{|B_0\setminus E|}{|B_0|}
 \leq a_{n,\beta}\rho_\theta(\alpha).
\label{wgt:inner-error}
\end{equation}
令 \(\omega=\Omega_{\mathcal D,w}(a_{n,\beta}\rho_\theta(\alpha))\)。若
\(w(B_0^\ast)=0\)，则 \(w(H)=0\)。否则由第一式 \eqref{wgt:outer-errors}，
\[
 w(B_0)\geq(1-\omega)w(B_0^\ast).
\]
当 \(\omega<1\) 时可知 \(w(B_0)>0\)，故由 \eqref{wgt:inner-error}，
\[
 w(E)\geq w(E\cap B_0)\geq(1-\omega)w(B_0)
 \geq(1-\omega)^2w(B_0^\ast).
\]
再由第二式 \eqref{wgt:outer-errors}，
\[
 w(H)\leq w(E)+w(B_0^\ast\setminus E)
 \leq\left(1+\frac{\omega}{(1-\omega)^2}\right)w(E).
\]
当 \(\omega\leq1/2\) 时，增量至多为 \(4\omega\)，证明
\eqref{wgt:Q-forward}--\eqref{wgt:R-forward}。

对中心方体，重新令 \(H=H_{\beta,\mathcal Q_c}(\alpha,E)\)，并选取其严格见证 \(B_0\)。见证聚拢结论给出 \(H\subset K\subset B_0\)，其中 \(K\) 与 \(B_0\) 同心，且
\[
 \frac{|K|}{|B_0|}\leq a_{n,\beta}\rho_\theta(\alpha)^n.
\]
若 \(w(B_0)=0\)，结论显然。否则由 \eqref{wgt:inner-error}，
\[
 w(E)\geq\bigl(1-\Omega_{\mathcal Q,w}(a_{n,\beta}\rho_\theta(\alpha))\bigr)w(B_0),
\]
而 \(w(H)\leq w(K)\) 给出 \eqref{wgt:Qc-forward}。
\end{proof}

\begin{theorem}[反向集中估计]\label{wgt:theorem-reverse}
设 \(w>0\) 几乎处处，且 \(\mathcal D\in\{\mathcal Q,\mathcal R\}\)。则每个 \(0<u<1\) 满足
\begin{equation}
 \Omega_{\mathcal D,w}(u)
 \leq1-\frac1{C^w_{\beta,\mathcal D}(1-u)}
 \leq C^w_{\beta,\mathcal D}(1-u)-1.
\label{wgt:reverse}
\end{equation}
若带权常数为 \(+\infty\)，约定右侧第一式中的倒数为 \(0\)。
\end{theorem}

\begin{proof}
取 \(D\in\mathcal D\)、\(F\subset D\) 且 \(|F|\leq u|D|\)，并令 \(E=D\setminus F\)、
\(\delta=|F|/|D|\)。由于 \(w>0\) 几乎处处，\(0<w(E)<\infty\)。对每个 \(x\in D\)，同一个
\(D\) 是合法的非中心见证，且
\[
 |D|^{\theta-1}|E|
 =(1-\delta)^{1-\theta}|E|^\theta>(1-u)|E|^\theta.
\]
严格性来自 \(0<1-\theta<1\)：若 \(\delta>0\)，第一处比较严格；若 \(\delta=0\)，则
\(u>0\) 使第二处比较严格。因此 \(D\subset H_{\beta,\mathcal D}(1-u,E)\)，从而
\[
 w(D)\leq C^w_{\beta,\mathcal D}(1-u)w(E).
\]
整理并对 \(D,F\) 取上确界即得第一式；第二式由 \(1-C^{-1}\leq C-1\) 得到。
\end{proof}

\begin{corollary}[权类刻画与指数保持]\label{wgt:corollary-Ainfty}
设 \(w>0\) 几乎处处。则
\[
 \lim_{\alpha\uparrow1}C^w_{\beta,\mathcal Q}(\alpha)=1
 \quad\Longleftrightarrow\quad w\in A_\infty,
\]
\[
 \lim_{\alpha\uparrow1}C^w_{\beta,\mathcal R}(\alpha)=1
 \quad\Longleftrightarrow\quad w\in A_\infty^\ast.
\]
此外，令
\[
 \Gamma_{\mathrm{conc}}(w,\mathcal D)
 :=\sup\{\gamma>0:\Omega_{\mathcal D,w}(u)=O(u^\gamma)\ (u\downarrow0)\},
\]
\[
 \Gamma_{\mathrm{halo}}(w,\mathcal D)
 :=\sup\{\gamma>0:C^w_{\beta,\mathcal D}(\alpha)-1
 =O((1-\alpha)^\gamma)\ (\alpha\uparrow1)\}.
\]
则对 \(\mathcal D\in\{\mathcal Q,\mathcal R\}\)，在扩展非负实数意义下，
\begin{equation}
 \Gamma_{\mathrm{halo}}(w,\mathcal D)=\Gamma_{\mathrm{conc}}(w,\mathcal D).
\label{wgt:exponent-preservation}
\end{equation}
\end{corollary}

\begin{proof}
通常的 \(A_\infty\) 条件和强 \(A_\infty^\ast\) 条件分别等价于相应集中模量在 \(0\) 处趋于零。充分性由定理 \ref{wgt:theorem-forward}，必要性由定理 \ref{wgt:theorem-reverse}。
同样，\eqref{wgt:reverse} 给出 \(\Gamma_{\mathrm{conc}}\geq\Gamma_{\mathrm{halo}}\)，而
\eqref{wgt:Q-forward}--\eqref{wgt:R-forward} 与 \(\rho_\theta(\alpha)\asymp_\theta1-\alpha\)
给出反向不等式。
\end{proof}

\begin{corollary}[定量集中条件]\label{wgt:corollary-power-conc}
设 \(A\geq1\)、\(0<\delta\leq1\)。若对所有可测的 \(F\subset Q\in\mathcal Q\) 且 \(w(Q)>0\) 有
\begin{equation}
 \frac{w(F)}{w(Q)}\leq A\left(\frac{|F|}{|Q|}\right)^\delta,
\label{wgt:power-concentration-cube}
\end{equation}
则 \(\alpha\) 充分接近 \(1\) 时
\[
 C^w_{\beta,\mathcal Q}(\alpha)-1\lesssim_{n,\beta,A,\delta}(1-\alpha)^\delta
\]
以及
\[
 C^w_{\beta,\mathcal Q_c}(\alpha)
 \lesssim_{n,\beta,A,\delta}(1-\alpha)^{n\delta}.
\]
若同一估计的 \(Q\) 全部替换为 \(R\)，并对所有可测的 \(F\subset R\in\mathcal R\) 且 \(w(R)>0\) 成立，则
\[
 C^w_{\beta,\mathcal R}(\alpha)-1\lesssim_{n,\beta,A,\delta}(1-\alpha)^\delta.
\]
矩形条件蕴含方体条件，因而在该更强假设下三个结论同时成立。
\end{corollary}

\begin{proof}
将相应的幂型集中界代入定理 \ref{wgt:theorem-forward} 即可。最后一个中心结论只是通用上界；其指数在特定齐次权下未必最佳。
\end{proof}

\begin{proposition}[Lebesgue 点局部化]\label{wgt:proposition-localization}
若 \(x_0\) 是 \(w\) 的正有限 Lebesgue 点，\(E,F\) 为有界可测集、\(0<|E|<\infty\)，且
\(F\subset H_{\beta,\mathcal B}(\alpha,E)\)，则对每个 \(r>0\)，
\begin{equation}
 x_0+rF\subset H_{\beta,\mathcal B}(\alpha,x_0+rE),
\qquad
 \frac{w(x_0+rF)}{w(x_0+rE)}\longrightarrow\frac{|F|}{|E|}\quad(r\downarrow0).
\label{wgt:localization}
\end{equation}
\end{proposition}

\begin{proof}
第一式由平移和伸缩不变性直接得到。取 \(R>0\) 使 \(E\cup F\subset Q(0,R)\)。对
\(A\in\{E,F\}\)，
\[
 \left|r^{-n}w(x_0+rA)-w(x_0)|A|\right|
 \leq r^{-n}\int_{Q(x_0,rR)}|w(y)-w(x_0)|\,dy\longrightarrow0.
\]
右端是固定常数乘以 Lebesgue 点平均振荡，故极限成立。由于 \(w(x_0)|E|>0\)，分母最终为正，遂可取比值。
\end{proof}

\begin{lemma}[径向奇性积分]\label{wgt:lemma-singular-integral}
设 \(d\geq1\)、\(0<a<d\)，并记 \(\omega_d=|B(0,1)|\)。对任意可测 \(A\subset\mathbb R^d\) 且 \(|A|<\infty\)，
\begin{equation}
 \int_A|x|^{-a}\,dx
 \leq\frac d{d-a}\omega_d^{a/d}|A|^{1-a/d}.
\label{wgt:singular-integral}
\end{equation}
\end{lemma}

\begin{proof}
令 \(m=|A|\)。层饼公式与 Tonelli 定理给出
\[
 \int_A|x|^{-a}\,dx
 \leq\int_0^\infty\min\{m,\omega_dt^{-d/a}\}\,dt.
\]
在 \(t_0=(\omega_d/m)^{a/d}\) 处分段积分，右侧恰为
\(d(d-a)^{-1}\omega_d^{a/d}m^{1-a/d}\)。
\end{proof}

\begin{proposition}[一坐标奇性权的非中心阶]\label{wgt:proposition-one-coordinate}
设 \(0<\gamma<1\)，\(w_\gamma(x)=|x_1|^{\gamma-1}\)。则
\begin{equation}
 C^{w_\gamma}_{\beta,\mathcal Q}(\alpha)-1\asymp_{n,\beta,\gamma}(1-\alpha)^\gamma,
 \qquad
 C^{w_\gamma}_{\beta,\mathcal R}(\alpha)-1\asymp_{n,\beta,\gamma}(1-\alpha)^\gamma.
\label{wgt:one-coordinate}
\end{equation}
\end{proposition}

\begin{proof}
记 \(v_\gamma(t)=|t|^{\gamma-1}\)。对任意区间 \(I\) 及 \(A\subset I\)，有
\begin{equation}
 \frac{v_\gamma(A)}{v_\gamma(I)}\leq C_\gamma
 \left(\frac{|A|}{|I|}\right)^\gamma.
\label{wgt:one-dimensional-concentration}
\end{equation}
事实上，若 \(d=\operatorname{dist}(0,I)\geq |I|\)，则 \(v_\gamma\) 在 \(I\) 上的上下界之比不超过
\(2^{1-\gamma}\)。其余情形中，若 \(I\) 跨过原点，则
\(v_\gamma(I)\geq|I|^\gamma/\gamma\)；若 \(I=(d,d+|I|)\) 且 \(0\leq d<|I|\)，
\((1+s)^\gamma-s^\gamma\) 在 \([0,1]\) 上递减，从而
\(v_\gamma(I)\geq(2^\gamma-1)|I|^\gamma/\gamma\)。再以引理
\ref{wgt:lemma-singular-integral} 的 \(d=1,a=1-\gamma\) 情形控制分子，即得
\eqref{wgt:one-dimensional-concentration}。

对任意轴平行矩形 \(R=I\times J\) 和 \(F\subset R\)，以层截函数
\(g(t)=|\{z\in J:(t,z)\in F\}|/|J|\) 表示 \(F\)。由层饼公式、
\eqref{wgt:one-dimensional-concentration} 以及 Jensen 不等式，
\begin{align*}
 \frac{w_\gamma(F)}{w_\gamma(R)}
 &=\int_0^1\frac{v_\gamma(\{g>s\})}{v_\gamma(I)}\,ds\\
 &\leq C_\gamma\int_0^1\left(\frac{|\{g>s\}|}{|I|}\right)^\gamma ds\\
 &\leq C_\gamma\left(\int_0^1\frac{|\{g>s\}|}{|I|}\,ds\right)^\gamma
 =C_\gamma\left(\frac{|F|}{|R|}\right)^\gamma.
\end{align*}
于是上界由推论 \ref{wgt:corollary-power-conc} 得到。

为得下界，令
\[
 E_h=(h,1+h)\times(0,1)^{n-1},\qquad
 F_h=(0,h)\times(0,1)^{n-1}.
\]
令 \(L_\alpha=\alpha^{-1/(n-\beta)}\)、
\(A_\alpha=\alpha^{-n/(n-\beta)}\)。若 \(h<L_\alpha-1\)，方体
\((0,1+h)^n\) 见证 \(F_h\subset H_{\beta,\mathcal Q}(\alpha,E_h)\)；若
\(h<A_\alpha-1\)，同一长方体见证
\(F_h\subset H_{\beta,\mathcal R}(\alpha,E_h)\)。另一方面，对每个 \(x\in E_h\)，集合
\(E_h\) 自身是两个非中心基的严格见证，因为其边长和体积均为 \(1\)，且
\(|E_h|^{\theta-1}|E_h|=1>\alpha\)。故在相应的两种条件下均有
\(E_h\cup F_h\subset H_{\beta,\mathcal B}(\alpha,E_h)\)。又
\[
 \frac{w_\gamma(E_h\cup F_h)}{w_\gamma(E_h)}-1
 =\frac{h^\gamma}{(1+h)^\gamma-h^\gamma}\asymp_\gamma h^\gamma
 \qquad(0<h\leq1/2).
\]
分别取 \(h=(L_\alpha-1)/2\) 及 \(h=(A_\alpha-1)/2\)，并用
\(L_\alpha-1\asymp A_\alpha-1\asymp1-\alpha\)，完成证明。
\end{proof}

\begin{proposition}[横向径向奇性权的中心阶]\label{wgt:proposition-radial}
写 \(x=(y,z)\in\mathbb R^k\times\mathbb R^{n-k}\)，其中 \(1\leq k\leq n\)，并设
\(0<a<k\)、\(w_{k,a}(y,z)=|y|^{-a}\)。则
\begin{equation}
 C^{w_{k,a}}_{\beta,\mathcal Q_c}(\alpha)
 \asymp_{n,\beta,k,a}(1-\alpha)^{n-a}.
\label{wgt:radial}
\end{equation}
\end{proposition}

\begin{proof}
设 \(Q=I\times J\) 是 \(\mathbb R^n\) 中边长为 \(\ell\) 的轴平行方体，\(F\subset Q\)。将
引理 \ref{wgt:lemma-singular-integral} 应用于 \(\mathbb R^k\) 中的截面。若
\(\operatorname{dist}(0,I)\geq2\sqrt k\ell\)，则 \(|y|^{-a}\) 在 \(I\) 上的上、下界之比由
\(k,a\) 控制，故
\[
 \frac{\int_A|y|^{-a}\,dy}{\int_I|y|^{-a}\,dy}
 \leq C_{k,a}\frac{|A|}{|I|}
 \leq C_{k,a}\left(\frac{|A|}{|I|}\right)^{1-a/k}
 \qquad(A\subset I).
\]
若 \(\operatorname{dist}(0,I)<2\sqrt k\ell\)，则 \(I\subset B(0,C_k\ell)\)，从而
\(|y|^{-a}\geq C_{k,a}\ell^{-a}\) 在 \(I\) 上成立，故
\[
 \int_I|y|^{-a}\,dy\geq c_{k,a}\ell^{k-a}.
\]
另一方面，引理 \ref{wgt:lemma-singular-integral} 给出
\[
 \int_A|y|^{-a}\,dy\leq C_{k,a}|A|^{1-a/k}.
\]
故两种情形均有同一集中不等式。对 \(F\) 的 \(z\)-截面作层饼，并对凹函数
\(r\mapsto r^{1-a/k}\) 使用 Jensen 不等式，遂得
\begin{equation}
 \frac{w_{k,a}(F)}{w_{k,a}(Q)}
 \leq C_{n,k,a}\left(\frac{|F|}{|Q|}\right)^{1-a/k}.
\label{wgt:radial-concentration}
\end{equation}
再设 \(sQ\) 与 \(Q\) 同心且边长为 \(s\ell\)。在远离奇点的情形，权在 \(Q\) 与
\(sQ\) 上均可比，故
\[
 \frac{w_{k,a}(sQ)}{w_{k,a}(Q)}\leq C_{k,a}s^n\leq C_{k,a}s^{n-a}.
\]
在近奇点情形，前述分母下界和引理 \ref{wgt:lemma-singular-integral} 分别给出
\[
 \int_I|y|^{-a}\,dy\geq c_{k,a}\ell^{k-a},\qquad
 \int_{I_s}|y|^{-a}\,dy\leq C_{k,a}(s\ell)^{k-a},
\]
其中 \(I_s\) 是 \(sQ\) 的 \(y\)-投影。再乘 \(z\)-方向的比例 \(s^{n-k}\)，得到
\begin{equation}
 \frac{w_{k,a}(sQ)}{w_{k,a}(Q)}\leq C_{n,k,a}s^{n-a},\qquad0<s<1.
\label{wgt:radial-dilation}
\end{equation}
这里 \eqref{wgt:radial-dilation} 是关键的额外输入。它不能由一般集中指数
\(1-a/k\) 代入推论 \ref{wgt:corollary-power-conc} 得到。

固定 \(E\)，令中心光晕非空，并取严格见证 \(B_0\)。由 \eqref{wgt:anchor-estimates} 和
\eqref{wgt:radial-concentration}，
\[
 \frac{w_{k,a}(B_0\setminus E)}{w_{k,a}(B_0)}
 \leq C_{n,k,a}\bigl(\rho_\theta(\alpha)\alpha^{-1/\theta}\bigr)^{1-a/k}=o(1).
\]
故 \(\alpha\) 充分接近 \(1\) 时，\(w_{k,a}(E)\geq w_{k,a}(B_0)/2\)。另一方面，
见证聚拢结论将整个光晕包含于 \(B_0\) 的同心缩放 \(s_\alpha B_0\)，其中
\(s_\alpha\leq C_{n,\beta}(1-\alpha)\)。于是由 \eqref{wgt:radial-dilation}，
\[
 \frac{w_{k,a}(H_{\beta,\mathcal Q_c}(\alpha,E))}{w_{k,a}(E)}
 \leq C_{n,\beta,k,a}(1-\alpha)^{n-a}.
\]
取 \(E=Q_0=[-1/2,1/2]^n\)。令 \(d_\alpha=1-\alpha^{1/\beta}\)。前文的内部见证构造给出
\(d_\alpha Q_0\subset H_{\beta,\mathcal Q_c}(\alpha,Q_0)\)（忽略零测边界）。由齐次性，
\[
 \frac{w_{k,a}(H_{\beta,\mathcal Q_c}(\alpha,Q_0))}{w_{k,a}(Q_0)}
 \geq d_\alpha^{n-a}\asymp_{\beta,k,a}(1-\alpha)^{n-a},
\]
从而完成证明。
\end{proof}

\section{多线性分数次光晕：对角约化与一般输入夹逼}
\label{sec:multilinear}

本节固定整数 $m\geq 1$ 和 $0<\beta<mn$，并写
\[
 \gamma:=\frac{\beta}{m}\in(0,n),\qquad
 q_{\beta,m}:=\frac{n}{mn-\beta}.
\]
令 $\mathcal B$ 为 $\mathcal Q,\mathcal R,\mathcal Q_c$ 中的任一基。所有积分和测度均指
Lebesgue 积分和 Lebesgue 测度；外层加权版本将在本节末尾单独说明。这里最重要的
归一化事实是：对角阈值是 $\alpha^m|E|^{\beta/n}$，而一般输入的自然尺度是各输入
测度的几何平均。二者不能互换。

\begin{definition}[多线性分数次极大算子]\label{ml:operator}
对局部可积函数 $f_1,\ldots,f_m$，定义
\[
 M_{\beta,m,\mathcal B}(f_1,\ldots,f_m)(x)
 :=\sup_{B\in\mathcal B(x)} |B|^{\beta/n-m}
      \prod_{i=1}^m\int_B |f_i(y)|\,dy .
\]
\end{definition}

\begin{lemma}[共同伸缩与自然顶层]\label{ml:scaling-top}
若 $r>0$，则对可测集 $E_1,\ldots,E_m$ 有
\[
 M_{\beta,m,\mathcal B}(\mathbf1_{rE_1},\ldots,\mathbf1_{rE_m})(rx)
 =r^\beta M_{\beta,m,\mathcal B}(\mathbf1_{E_1},\ldots,\mathbf1_{E_m})(x).
\]
若 $0<|E_i|<\infty$，则对每个 $x$，
\[
 M_{\beta,m,\mathcal B}(\mathbf1_{E_1},\ldots,\mathbf1_{E_m})(x)
 \leq \prod_{i=1}^m |E_i|^{\beta/(mn)}.
\]
\end{lemma}

\begin{proof}
把 $rx$ 的见证写成 $rB$。每个积分给出因子 $r^n$，而
$|rB|^{\beta/n-m}=r^{\beta-nm}|B|^{\beta/n-m}$，故得到第一个等式；反向缩放给出等号。
另一方面，对每个 $i$，线性自然顶层估计（参数为 $\gamma$）给出
\[
 |B|^{\gamma/n-1}|E_i\cap B|\leq |E_i|^{\gamma/n}.
\]
相乘后再取上确界即得第二个结论。
\end{proof}

\subsection{对角资料}

\begin{definition}[对角光晕]\label{ml:diagonal-halo}
对 $0<\alpha<1$ 及 $0<|E|<\infty$，令
\[
 H^{\Delta}_{\beta,m,\mathcal B}(\alpha,E)
 :=\left\{x:
 M_{\beta,m,\mathcal B}(\mathbf1_E,\ldots,\mathbf1_E)(x)
 >\alpha^m|E|^{\beta/n}\right\},
\]
并令
\[
 C^{\Delta}_{\beta,m,\mathcal B}(\alpha)
 :=\sup_{0<|E|<\infty}
 \frac{|H^{\Delta}_{\beta,m,\mathcal B}(\alpha,E)|}{|E|}.
\]
\end{definition}

\begin{theorem}[精确对角约化]\label{ml:diagonal-reduction}
对每个 $f\in L^1_{\mathrm{loc}}(\mathbb R^n)$，
\[
 M_{\beta,m,\mathcal B}(f,\ldots,f)
 =\bigl(M_{\gamma,\mathcal B}f\bigr)^m.
\]
因而对所有 $0<\alpha<1$ 及 $0<|E|<\infty$，
\[
 H^{\Delta}_{\beta,m,\mathcal B}(\alpha,E)
 =H_{\gamma,\mathcal B}(\alpha,E),
 \qquad
 C^{\Delta}_{\beta,m,\mathcal B}(\alpha)
 =C_{\gamma,\mathcal B}(\alpha).
\]
\end{theorem}

\begin{proof}
因为非负数的 $m$ 次幂保持上确界，
\[
 \begin{split}
 M_{\beta,m,\mathcal B}(f,\ldots,f)(x)
 &=\sup_{B\in\mathcal B(x)}
 \left(|B|^{\beta/(mn)-1}\int_B|f|\right)^m\\
 &=\bigl(M_{\gamma,\mathcal B}f(x)\bigr)^m.
 \end{split}
\]
又 $\alpha^m|E|^{\beta/n}=(\alpha|E|^{\gamma/n})^m$，且 $s\mapsto s^m$
严格递增，故两个严格水平集相同。最后对 $E$ 取上确界即可。
\end{proof}

\begin{corollary}[对角模型的端点结论]\label{ml:diagonal-endpoint}
当 $\alpha\uparrow1$ 时，
\[
 C^{\Delta}_{\beta,m,\mathcal Q}(\alpha)-1\asymp 1-\alpha,
 \quad
 C^{\Delta}_{\beta,m,\mathcal R}(\alpha)-1\asymp 1-\alpha,
 \quad
 C^{\Delta}_{\beta,m,\mathcal Q_c}(\alpha)\asymp(1-\alpha)^n,
\]
其中隐含常数只依赖于 $n,\beta,m$。在 $n=1$ 且
$2^{-\beta/m}<\alpha<1$ 时，
\[
 C^{\Delta}_{\beta,m,\mathcal Q}(\alpha)
 =C^{\Delta}_{\beta,m,\mathcal R}(\alpha)
 =2\alpha^{-m/(m-\beta)}-1,
\]
以及
\[
 C^{\Delta}_{\beta,m,\mathcal Q_c}(\alpha)
 =\max\left\{1-\alpha^{m/\beta},\
                 \alpha^{-m/(m-\beta)}-1\right\}.
\]
\end{corollary}

\begin{proof}
这是定理~\ref{ml:diagonal-reduction} 对参数 $\gamma=\beta/m$ 的直接应用。
一维公式同样由线性一维闭式代入 $\gamma$ 得到。
\end{proof}

若 $Q_0=[-1/2,1/2]^n$，令
\[
 c_{\beta,m}:=\min\{\beta/m,n-\beta/m\},\qquad
 k_{\beta,m}:=\lfloor2c_{\beta,m}\rfloor,
\]
其中 $2c_{\beta,m}=n$ 时约定 $k_{\beta,m}=n$。固定 $Q_0$ 的线性切向极限遂给出
\[
 \lim_{\alpha\uparrow1}
 \frac{|H^{\Delta}_{\beta,m,\mathcal Q_c}(\alpha,Q_0)|}{(1-\alpha)^n}
 =\frac{2^n}{k_{\beta,m}!(2c_{\beta,m})^{n-k_{\beta,m}}}.
\]
这是固定资料的极限；它只给出全体对角常数的下极限下界，不能推出全体集合的
尖锐首项等式。

\subsection{一般输入及其自然的几何平均尺度}

\begin{definition}[几何平均归一化]\label{ml:geometric-mean}
对 $0<|E_i|<\infty$，令
\[
 G(E_1,\ldots,E_m):=\prod_{i=1}^m|E_i|^{1/m},
\]
并设
\[
 \begin{split}
 H^{\mathrm{gm}}_{\beta,m,\mathcal B}(\alpha;E_1,\ldots,E_m)
 :=\biggl\{x:\;&M_{\beta,m,\mathcal B}(\mathbf1_{E_1},\ldots,\mathbf1_{E_m})(x)\\
 &>\alpha^m\prod_{i=1}^m|E_i|^{\beta/(mn)}\biggr\},
 \end{split}
\]
\[
 C^{\mathrm{gm}}_{\beta,m,\mathcal B}(\alpha)
 :=\sup_{0<|E_i|<\infty}
 \frac{|H^{\mathrm{gm}}_{\beta,m,\mathcal B}(\alpha;E_1,\ldots,E_m)|}
 {G(E_1,\ldots,E_m)}.
\]
\end{definition}

\begin{lemma}[尺度不变性]\label{ml:gm-scaling}
对 $r>0$，
\begin{align*}
 H^{\mathrm{gm}}_{\beta,m,\mathcal B}(\alpha;rE_1,\ldots,rE_m)
 &=rH^{\mathrm{gm}}_{\beta,m,\mathcal B}(\alpha;E_1,\ldots,E_m),\\
 G(rE_1,\ldots,rE_m)&=r^nG(E_1,\ldots,E_m).
\end{align*}
\end{lemma}

\begin{proof}
引理~\ref{ml:scaling-top} 表明算子与阈值同乘 $r^\beta$；第二式是测度伸缩公式。
\end{proof}

\begin{theorem}[一般输入的精确夹逼]\label{ml:general-sandwich}
对所有 $0<\alpha<1$，
\[
 C_{\gamma,\mathcal B}(\alpha)
 \leq C^{\mathrm{gm}}_{\beta,m,\mathcal B}(\alpha)
 \leq C_{\gamma,\mathcal B}(\alpha^m).
\]
更精确地，任意 $E_1,\ldots,E_m$ 满足
\[
 H^{\mathrm{gm}}_{\beta,m,\mathcal B}(\alpha;E_1,\ldots,E_m)
 \subseteq\bigcap_{i=1}^m H_{\gamma,\mathcal B}(\alpha^m,E_i).
\]
\end{theorem}

\begin{proof}
固定 $B\in\mathcal B(x)$，并令
\[
 u_i(B):=\frac{|B|^{\gamma/n-1}|E_i\cap B|}{|E_i|^{\gamma/n}}.
\]
由引理~\ref{ml:scaling-top}，$0\leq u_i(B)\leq1$。多线性严格见证条件恰为
$\prod_i u_i(B)>\alpha^m$。若某个 $u_i(B)\leq\alpha^m$，则该乘积不超过
$\alpha^m$，矛盾。因此每个 $u_i(B)>\alpha^m$，而同一个 $B$ 同时见证右端所有
线性光晕，证明集合包含。

由此
\[
 |H^{\mathrm{gm}}_{\beta,m,\mathcal B}|
 \leq\min_i|H_{\gamma,\mathcal B}(\alpha^m,E_i)|
 \leq C_{\gamma,\mathcal B}(\alpha^m)\min_i|E_i|
 \leq C_{\gamma,\mathcal B}(\alpha^m)G(E_1,\ldots,E_m).
\]
取上确界得到上界。令 $E_1=\cdots=E_m=E$，定理~\ref{ml:diagonal-reduction}
给出下界。
\end{proof}

\begin{corollary}[一般输入的端点指数]\label{ml:general-endpoint}
当 $\alpha\uparrow1$ 时，
\[
 C^{\mathrm{gm}}_{\beta,m,\mathcal Q}(\alpha)-1\asymp1-\alpha,
 \qquad
 C^{\mathrm{gm}}_{\beta,m,\mathcal R}(\alpha)-1\asymp1-\alpha,
\]
\[
 C^{\mathrm{gm}}_{\beta,m,\mathcal Q_c}(\alpha)
 \asymp(1-\alpha)^n.
\]
\end{corollary}

\begin{proof}
由定理~\ref{ml:general-sandwich}，并注意
$1-\alpha^m=(1-\alpha)(1+\cdots+\alpha^{m-1})\asymp_m1-\alpha$。
\end{proof}

\subsection{外层权与高矩的对角传递}

\begin{definition}[外层加权多线性常数]\label{ml:weighted-definition}
令 $w\geq0$ 局部可积且不恒为零。若 $0<w(E)<\infty$，定义
\[
 C^{w,\Delta}_{\beta,m,\mathcal B}(\alpha)
 :=\sup_{\substack{0<|E|<\infty\\0<w(E)<\infty}}
 \frac{w(H^{\Delta}_{\beta,m,\mathcal B}(\alpha,E))}{w(E)}.
\]
对一般输入，令
\[
 C^{w,\mathrm{gm}}_{\beta,m,\mathcal B}(\alpha)
 :=\sup_{\substack{0<|E_i|<\infty\\0<w(E_i)<\infty}}
 \frac{w(H^{\mathrm{gm}}_{\beta,m,\mathcal B}(\alpha;E_1,\ldots,E_m))}
 {\prod_{i=1}^m w(E_i)^{1/m}}.
\]
这些定义只在光晕仍由 Lebesgue 分数次平均定义的\emph{外层权}意义下使用。
\end{definition}

\begin{theorem}[外层权的对角约化与夹逼]\label{ml:weighted-sandwich}
对 $0<\alpha<1$，
\[
 C^{w,\Delta}_{\beta,m,\mathcal B}(\alpha)=C^w_{\gamma,\mathcal B}(\alpha),
 \qquad
 C^w_{\gamma,\mathcal B}(\alpha)
 \leq C^{w,\mathrm{gm}}_{\beta,m,\mathcal B}(\alpha)
 \leq C^w_{\gamma,\mathcal B}(\alpha^m).
\]
\end{theorem}

\begin{proof}
第一个恒等式来自定理~\ref{ml:diagonal-reduction}。对第二个式子，定理~\ref{ml:general-sandwich}
的集合包含和单调性给出
\[
 w(H^{\mathrm{gm}})\leq\min_iw(H_{\gamma,\mathcal B}(\alpha^m,E_i))
 \leq C^w_{\gamma,\mathcal B}(\alpha^m)\min_iw(E_i)
 \leq C^w_{\gamma,\mathcal B}(\alpha^m)\prod_iw(E_i)^{1/m}.
\]
对角输入给出反向不等式。
\end{proof}

\begin{corollary}[外层权的端点刻画与指数]\label{ml:weighted-endpoint}
设 $w>0$ 几乎处处，并设线性外层权结论已经在相同的 Lebesgue 平均模型下成立。
则
\[
 \lim_{\alpha\uparrow1}C^{w,\Delta}_{\beta,m,\mathcal Q}(\alpha)=1
 \quad\Longleftrightarrow\quad w\in A_\infty,
\]
\[
 \lim_{\alpha\uparrow1}C^{w,\Delta}_{\beta,m,\mathcal R}(\alpha)=1
 \quad\Longleftrightarrow\quad w\in A_\infty^*.
\]
将 $C^{w,\Delta}$ 换成 $C^{w,\mathrm{gm}}$ 后这两个等价仍成立。若对
$\mathcal D\in\{\mathcal Q,\mathcal R\}$ 定义
\[
 \Gamma^{w,\mathrm{gm}}_{\mathcal D}:=
 \sup\{\delta>0:C^{w,\mathrm{gm}}_{\beta,m,\mathcal D}(\alpha)-1
 =O((1-\alpha)^\delta)\},
\]
则 $\Gamma^{w,\mathrm{gm}}_{\mathcal D}$ 等于对应线性外层权模型的集中指数；
对角模型亦然。
\end{corollary}

\begin{proof}
定理~\ref{ml:weighted-sandwich} 把对角模型完全化为参数 $\gamma$ 的线性模型，并把
一般输入夹在其参数 $\alpha$ 和 $\alpha^m$ 的两个值之间。由于
$1-\alpha^m\asymp_m1-\alpha$，线性外层权刻画和指数结论同时传递。
\end{proof}

这里的结论不涉及以 $w(x)dx$ 取平均的内禀加权极大算子。

为便于陈述高矩结论，令
\[
 \mathfrak F^{\Delta}_{\beta,m,\mathcal B}(p)
 :=\sup_{0<|E|<\infty}\frac1{|E|}
 \int_{\mathbb R^n}
 \left(\frac{M_{\beta,m,\mathcal B}(\mathbf1_E,\ldots,\mathbf1_E)(x)}
 {|E|^{\beta/n}}\right)^p dx .
\]
\begin{proposition}[对角高矩恒等式]\label{ml:diagonal-moments}
当 $p>q_{\beta,m}$ 时，
\[
 \mathfrak F^{\Delta}_{\beta,m,\mathcal B}(p)
 =\mathfrak F_{\gamma,\mathcal B}(mp),
\]
其中右端是参数 $\gamma$ 的线性归一化高矩。由定理~\ref{mom:fixed-cube-moment}
（以参数 $\gamma$ 应用）无条件得到
\[
 \lim_{p\to\infty}p^n\int_{\mathbb R^n}
 M_{\beta,m,\mathcal Q_c}(\mathbf1_{Q_0},\ldots,\mathbf1_{Q_0})(x)^p\,dx
 =\frac{n!}{m^n}\,
 \frac{2^n}{k_{\beta,m}!(2c_{\beta,m})^{n-k_{\beta,m}}}.
\]
\end{proposition}

\begin{proof}
由定理~\ref{ml:diagonal-reduction} 和
$|E|^{\beta/n}=(|E|^{\gamma/n})^m$，被积函数逐点等于线性归一化函数的 $mp$ 次幂。
积分后取上确界即得第一式。第二式对 $q=mp$ 应用固定 $Q_0$ 的线性高矩首项，并用
$q^n=m^np^n$。
\end{proof}

\subsection{公共交集归一化的退化}

\begin{proposition}[正分数阶下公共交集分母必然发散]\label{ml:intersection-divergence}
设 $m\geq2$、$\beta>0$，并令 $\mathcal B\in\{\mathcal Q,\mathcal R,\mathcal Q_c\}$。
对任意 $\Lambda>0$，定义
\[
 \mathcal C^{\cap}_{\beta,m,\mathcal B}(\Lambda)
 :=\sup_{\substack{0<|\cap_iE_i|<\infty\\0<|E_i|<\infty}}
 \frac{\bigl|\{M_{\beta,m,\mathcal B}(\mathbf1_{E_1},\ldots,\mathbf1_{E_m})
 >\Lambda|\cap_iE_i|^{\beta/n}\}\bigr|}{|\cap_iE_i|}.
\]
则 $\mathcal C^{\cap}_{\beta,m,\mathcal B}(\Lambda)=\infty$。
\end{proposition}

\begin{proof}
以第一坐标的小数部分把 $\mathbb R^n$ 周期地分成 $m$ 个集合
\[
 A_i^0:=\{x:\{x_1\}\in[(i-1)/m,i/m)\},\qquad 1\leq i\leq m.
\]
若轴平行方体 $Q$ 的边长为整数，则 $|A_i^0\cap Q|=|Q|/m$；此结论对平移的
方体也成立。令 $F=[-1/2,1/2]^n$，并对正整数 $N$ 令
\[
 D_N=[-2N,2N]^n,\qquad E_i^{(N)}:=D_N\setminus(A_i^0\setminus F).
\]
则每个 $E_i^{(N)}$ 有限测度且 $\bigcap_iE_i^{(N)}=F$。

对 $x\in[-N/2,N/2]^n$，取以 $x$ 为中心、边长 $2N$ 的方体 $Q_x$。它包含于 $D_N$ 并
包含 $F$，故它是三种基的合法见证，且
\[
 |E_i^{(N)}\cap Q_x|\geq(1-1/m)|Q_x|.
\]
于是
\[
 M_{\beta,m,\mathcal B}(\mathbf1_{E_1^{(N)}},\ldots,\mathbf1_{E_m^{(N)}})(x)
 \geq (1-1/m)^m(2N)^\beta.
\]
右端随 $N$ 发散。因此对充分大的 $N$，整个 $[-N/2,N/2]^n$ 都落入定义中的水平集；
该水平集商至少为 $N^n$，结论成立。
\end{proof}

\begin{remark}\label{ml:zero-order-boundary}
上面证明使用的是 $(2N)^\beta$。在 $\beta=0$ 时，这个增长因子消失，故该构造本身
并不说明零阶公共交集模型是否发散；这正是正分数阶与零阶之间的结构分界。
\end{remark}

\section{弱型控制、全体高矩首项与限制型范数}\label{mom:section}

本节的作用是把光晕的顶层行为转化为归一化示性函数的高矩行为。这里的弱型估计只针对 Lebesgue 测度；它不是关于任意底测度的一致结论。记
\[
 q_\beta:=\frac{n}{n-\beta}.
\]
对一个几何基 $\mathcal B$ 和 $p>q_\beta$，定义限制型泛函
\begin{equation}
 \mathfrak F_{\beta,\mathcal B}(p):=
 \sup_{0<|E|<\infty}\frac1{|E|}
 \int_{\mathbb R^n}\left(\frac{M_{\beta,\mathcal B}\mathbf 1_E(x)}{|E|^{\beta/n}}\right)^p\,dx .
 \label{mom:functional}
\end{equation}
若 $r_p^{-1}=p^{-1}+\beta/n$，则 $1<r_p<n/\beta$，且 $\mathfrak F_{\beta,\mathcal B}(p)^{1/p}$ 正是此 Sobolev 线上限制于示性函数的 $L^{r_p}\to L^p$ 最优常数。

\begin{lemma}[有限中心方体的选择]\label{mom:selection}
设 $Q_j=x_j+[-r_j,r_j]^n$（$1\le j\le N$）为有限个中心方体。存在子族 $\{Q_{j_\nu}\}_{\nu=1}^M$，满足：
\begin{enumerate}
\item 对每个 $j$，存在 $\nu$ 使 $r_{j_\nu}\ge r_j$ 且 $x_j\in Q_{j_\nu}$；
\item $\sum_{\nu=1}^M\mathbf 1_{Q_{j_\nu}}\le2^n$ 几乎处处。
\end{enumerate}
\end{lemma}
\begin{proof}
按半边长递减排列。依次扫描方体；其中心不属于此前已选方体时才保留它。被舍弃的方体即满足第一项。

取一个不在所有坐标超平面 $\{y_k=(x_j)_k\}$ 上的点 $y$。对于包含 $y$ 的已选方体，按 $x_j-y$ 的坐标符号分为 $2^n$ 个正交象限。若同一象限内 $Q_i,Q_j$ 均包含 $y$ 且 $r_i\ge r_j$，则逐坐标有
\[
 |(x_i)_k-(x_j)_k|\le r_i,
\]
所以 $x_j\in Q_i$。在扫描中 $Q_j$ 本应被舍弃，矛盾。因此每个象限至多对应一个已选方体。除去零测的坐标超平面并得第二项。
\end{proof}

\begin{lemma}[紧集上的有限稳定见证]\label{mom:stable-witness}
设 $f\in L^1(\mathbb R^n)$，$\lambda>0$，且
\[
 \Omega_\lambda=\{x:M_{\beta,\mathcal Q_c}f(x)>\lambda\}.
\]
对每个紧集 $K\subset\Omega_\lambda$ 和 $0<\varepsilon<1$，存在有限个 $x_j\in K$、$r_j>0$ 及开集 $U_j\ni x_j$，使 $K\subset\bigcup_jU_j$，$U_j\subset x_j+(-\varepsilon r_j,\varepsilon r_j)^n$，且对每个 $z\in U_j$，方体 $z+[-r_j,r_j]^n$ 都严格见证 $M_{\beta,\mathcal Q_c}f(z)>\lambda$。
\end{lemma}
\begin{proof}
对每个 $x\in K$ 选择一个严格见证半边长 $r_x$。固定半边长的卷积平均关于中心连续，故该严格不等式在 $x$ 的某个开邻域内仍成立。再与 $x+(-\varepsilon r_x,\varepsilon r_x)^n$ 相交，最后对 $K$ 取有限子覆盖。
\end{proof}

\begin{theorem}[中心方体的分数弱型估计]\label{mom:centered-weak}
对 $f\in L^1(\mathbb R^n)$ 和 $\lambda>0$，
\begin{equation}
 \lambda\,|\{M_{\beta,\mathcal Q_c}f>\lambda\}|^{1/q_\beta}
 \le 2^{n/q_\beta}\|f\|_1=2^{n-\beta}\|f\|_1.
 \label{mom:weak}
\end{equation}
因此
\begin{equation}
 C_{\beta,\mathcal Q_c}(\alpha)\le 2^n\alpha^{-q_\beta}\qquad(0<\alpha<1).
 \label{mom:centered-tauberian-bound}
\end{equation}
\end{theorem}
\begin{proof}
令 $K$ 为严格水平集的紧子集，并应用引理~\ref{mom:stable-witness}。对所得有限见证族使用引理~\ref{mom:selection}。若 $z\in K$，则某个稳定见证的中心落在某个已选方体中，且其半边长不超过已选方体的半边长；所以
\[
 K\subset\bigcup_{\nu=1}^M\bigl(x_{j_\nu}+[-(1+\varepsilon)r_{j_\nu},(1+\varepsilon)r_{j_\nu}]^n\bigr).
\]
记 $a_\nu=\int_{Q_{j_\nu}}|f|$。严格见证给出
\[
 \lambda |Q_{j_\nu}|^{1/q_\beta}<a_\nu.
\]
于是由 $q_\beta>1$、Minkowski 不等式和引理~\ref{mom:selection}，
\[
 \begin{split}
 \lambda |K|^{1/q_\beta}
 &\le(1+\varepsilon)^{n/q_\beta}
 \Bigl(\sum_\nu a_\nu^{q_\beta}\Bigr)^{1/q_\beta}\\
 &\le(1+\varepsilon)^{n/q_\beta}
 \int |f(y)|\Bigl(\sum_\nu\mathbf1_{Q_{j_\nu}}(y)\Bigr)^{1/q_\beta}\,dy\\
 &\le(1+\varepsilon)^{n/q_\beta}2^{n/q_\beta}\|f\|_1.
 \end{split}
\]
令 $\varepsilon\downarrow0$，再用开集的 Lebesgue 内正则性，得到 \eqref{mom:weak}。令 $f=\mathbf1_E$、$\lambda=\alpha|E|^{\beta/n}$，并利用 $q_\beta\beta/n=q_\beta-1$，即得 \eqref{mom:centered-tauberian-bound}。
\end{proof}

\begin{remark}\label{mom:weak-scope}
常数 $2^{n-\beta}$ 来自有限中心方体的 $2^n$ 重叠界。定理没有断言该常数是 Lebesgue Tauberian 常数或任意测度弱型常数的最佳值。
\end{remark}

\begin{proposition}[全域低阈值支配]\label{mom:low-level}
存在仅依赖 $n$ 的常数 $K_n$，使
\[
 C_{\beta,\mathcal Q}(\alpha)\le3^n\alpha^{-q_\beta},\qquad
 C_{\beta,\mathcal Q_c}(\alpha)\le2^n\alpha^{-q_\beta},
\]
并且
\[
 C_{\beta,\mathcal R}(\alpha)
 \le K_n^2\alpha^{-2/(1-\beta/n)}.
\]
\end{proposition}
\begin{proof}
中心情形已经由定理~\ref{mom:centered-weak} 给出。以下证明非中心方体估计，以便记录低阈值支配实际使用的选择步骤。令
\[
 \Omega=H_{\beta,\mathcal Q}(\alpha,E),\qquad
 \lambda=\alpha|E|^{\beta/n}.
\]
严格见证在小的外扩下仍严格，故对每个紧集 $K\subset\Omega$ 可取有限个严格见证方体 $Q_1,\ldots,Q_N$，使 $K\subset\bigcup_j\operatorname{int}Q_j$。按边长递减贪心地保留与此前已选方体不相交的方体，记所得子族为 $Q_{j_1},\ldots,Q_{j_M}$。每个被舍弃的 $Q_j$ 与某个不小于它的已选方体相交，故逐坐标比较区间得到 $Q_j\subset3Q_{j_\nu}$。于是
\[
 |K|\le3^n\sum_{\nu=1}^M|Q_{j_\nu}|.
\]
令 $s_\nu=|E\cap Q_{j_\nu}|$。严格见证和 $q_\beta=(1-\beta/n)^{-1}$ 给出
\[
 |Q_{j_\nu}|<\lambda^{-q_\beta}s_\nu^{q_\beta}.
\]
所选方体两两不交，故 $\sum_\nu s_\nu\le|E|$；又 $q_\beta>1$，从而
\[
 \lambda^{-q_\beta}\sum_\nu s_\nu^{q_\beta}
 \le\lambda^{-q_\beta}|E|^{q_\beta}
 =\alpha^{-q_\beta}|E|.
\]
对紧 $K$ 取上确界即得第一个不等式。

对矩形，若 $s=|E\cap R|$，则
\[
 |R|^{\beta/n-1}s/|E|^{\beta/n}
 \le(s/|R|)^{1-\beta/n}.
\]
故相应光晕包含于经典强极大函数的水平集 $\{M_{0,\mathcal R}\mathbf1_E>\alpha^{1/(1-\beta/n)}\}$。经典强极大函数在 $L^2$ 上有仅依赖 $n$ 的范数界；这可由一维 Hardy--Littlewood 极大函数逐坐标迭代得到，亦见 de Guzm\'an \cite{deGuzman1975}。再用 Chebyshev 不等式即得结论。
\end{proof}

\begin{theorem}[固定方体的高矩首项]\label{mom:fixed-cube-moment}
设 $Q_0=[-1/2,1/2]^n$。由定理~\ref{tan:unit-cube-limit}，
\[
 |H_{\beta,\mathcal Q_c}(1-t,Q_0)|=|S_{n,c}|t^n+o(t^n)\quad(t\downarrow0),
\]
其中 $c=\min\{\beta,n-\beta\}$，$k=\lfloor2c\rfloor$，并且
\begin{equation}
 \lim_{p\to\infty}p^n\int_{\mathbb R^n}
 \bigl(M_{\beta,\mathcal Q_c}\mathbf1_{Q_0}\bigr)^p\,dx
 =n!|S_{n,c}|=
 \frac{n!\,2^n}{k!(2c)^{n-k}}.
 \label{mom:fixed-cube-asymptotic}
\end{equation}
\end{theorem}
\begin{proof}
层析公式和 $|Q_0|=1$ 给出
\[
 \int(M_{\beta,\mathcal Q_c}\mathbf1_{Q_0})^p
 =p\int_0^1\alpha^{p-1}|H_{\beta,\mathcal Q_c}(\alpha,Q_0)|\,d\alpha.
\]
在充分靠近 $1$ 的区间，中心端点上界给出 $|H_{\beta,\mathcal Q_c}(\alpha,Q_0)|\le A(1-\alpha)^n$；远离 $1$ 的部分由命题~\ref{mom:low-level} 指数衰减。故可在近端使用支配收敛，再用
\[
 p\int_0^1\alpha^{p-1}(1-\alpha)^n\,d\alpha=pB(p,n+1)\sim n!p^{-n}.
\]
这证明结论。
\end{proof}

\begin{theorem}[限制型高矩渐近]\label{mom:restricted-growth}
记 $L_{n,\beta}$ 如定理~\ref{gl:main}。当 $p\to\infty$ 时，
\[
 \mathfrak F_{\beta,\mathcal Q}(p)=1+O_{n,\beta}(p^{-1}),\qquad
 \mathfrak F_{\beta,\mathcal R}(p)=1+O_{n,\beta}(p^{-1}),
\]
且
\[
 \mathfrak F_{\beta,\mathcal Q_c}(p)
 =n!L_{n,\beta}p^{-n}+O_{n,\beta}(p^{-n-1}).
\]
中心公式的误差还在任意 $\beta\in J\Subset(0,n)$ 上一致（$p$ 充分大取决于 $n,J$）。
此外，非中心方体和强矩形满足
\[
 \liminf_{p\to\infty}p\bigl(\mathfrak F_{\beta,\mathcal Q}(p)-1\bigr)
 \ge\frac{2n}{n-\beta},
 \qquad
 \liminf_{p\to\infty}p\bigl(\mathfrak F_{\beta,\mathcal R}(p)-1\bigr)
 \ge\frac{2n^2}{n-\beta}.
\]
在 $n=1$，由一维精确公式，两个相应极限分别为
\[
 \lim_{p\to\infty}p\bigl(\mathfrak F_{\beta,\mathcal Q}(p)-1\bigr)=\frac{2}{1-\beta},
 \qquad
 \lim_{p\to\infty}p\mathfrak F_{\beta,\mathcal Q_c}(p)
 =\max\!\left\{\frac1\beta,\frac1{1-\beta}\right\}.
\]
\end{theorem}
\begin{proof}
对任意 $E$，自然顶层和层析公式给出
\[
 \frac1{|E|}\int\left(\frac{M_{\beta,\mathcal B}\mathbf1_E}{|E|^{\beta/n}}\right)^p
 =p\int_0^1\alpha^{p-1}
 \frac{|H_{\beta,\mathcal B}(\alpha,E)|}{|E|}\,d\alpha.
\]
对上确界使用端点光晕上界，并由命题~\ref{mom:low-level} 支配远离 $1$ 的区间。Beta 积分
\[
 p\int_0^1\alpha^{p-1}(1-\alpha)^j\,d\alpha\sim j!p^{-j}
\]
给出非中心的两个上界。中心情形使用定理~\ref{gl:main}：对每个输入，
近端层析积分至多
$L_{n,\beta}pB(p,n+1)+O_{n,J}(pB(p,n+2))$，
远端由弱型界一致指数衰减。下界取固定 $Q_0$ 并用其 $N=1$ 一致截断，得到同样的表达式。
由于 $pB(p,n+1)=n!\prod_{j=1}^n(p+j)^{-1}$、$pB(p,n+2)=O_n(p^{-n-1})$，
中心结论及紧参数一致性成立。这里仅将上确界放在积分外作上界，没有交换二者为等式。
为得到两个非中心下界，固定 $E=[0,1]^n$。实际测试方体包含给出
\[
 C_{\beta,\mathcal Q}(1-t)\ge
 \bigl(2(1-t)^{-1/(n-\beta)}-1\bigr)^n
 =1+\frac{2n}{n-\beta}t+O(t^2),
\]
以及
\[
 C_{\beta,\mathcal R}(1-t)\ge
 1+2n\bigl((1-t)^{-n/(n-\beta)}-1\bigr)
 =1+\frac{2n^2}{n-\beta}t+O(t^2).
\]
把这两个固定测试集下界代入层析公式，并使用
\[
 p\int_0^1\alpha^{p-1}(1-\alpha)\,d\alpha=\frac1{p+1},
\]
即得所列两个 \(\liminf\)。一维结论由一维精确常数在固定近端区间内达到，再作同一 Beta 积分得到。
\end{proof}

\begin{remark}\label{mom:global-gap}
定理~\ref{gl:main}闭合中心全体首项后，限制型泛函也具有匹配的全体首项。
下一系数仍不能仅由固定输入 $Q_0$ 的高矩展开推出，因为取上确界的输入可以随 $p$ 改变。
\end{remark}

\begin{corollary}[二维半阶的全体高矩次阶夹逼]\label{mom:pair-bracket}
当 $n=2$、$\beta=1$ 时，
\[
 0\le\liminf_{p\to\infty}p^3[\mathfrak F_{1,\mathcal Q_c}(p)-4p^{-2}]
 \le\limsup_{p\to\infty}p^3[\mathfrak F_{1,\mathcal Q_c}(p)-4p^{-2}]\le8.
\]
\end{corollary}
\begin{proof}
定理~\ref{pb:volume}对每个输入给出统一近端层析上界。由
$\mathcal U(1-t)=2t^2+(10/3)t^3+O(t^4)$ 及低水平弱型支配，
\[
 \mathfrak F_{1,\mathcal Q_c}(p)
 \le2pB(p,3)+\tfrac{10}{3}pB(p,4)+O(p^{-4})
 =4p^{-2}+8p^{-3}+O(p^{-4}).
\]
下界使用固定 $Q_0$ 的面积展开 $2t^2+2t^3+O(t^4)$，
而非把矩形上确界放进积分；它给出
$\mathfrak F_{1,\mathcal Q_c}(p)\ge2pB(p,3)+2pB(p,4)+O(p^{-4})
=4p^{-2}+O(p^{-4})$。两个不等式合并即得结论。
\end{proof}


\begin{corollary}[高矩层析包络的显式改进]\label{mom:sixth-envelope}
固定 $n=2$、$\beta=1$，令 $a_0=1/\sqrt2$。对 $p>2$ 定义
\[
 \mathcal B(p)=\frac{4p}{p-2}a_0^{p-2}
 +p\int_{a_0}^1\alpha^{p-1}\mathcal U(\alpha)\,d\alpha.
\]
则
\[
 \mathfrak F_{1,\mathcal Q_c}(p)\le\mathcal B(p)-\mathcal D(p),\qquad
 \mathcal D(p)=\frac p{768}\int_{99/100}^1\alpha^{p-1}(1-\alpha)^6\,d\alpha>0,
\]
且
\[
 \mathcal D(p)=\frac{15}{16}\prod_{j=1}^6(p+j)^{-1}
 +O\bigl((99/100)^p\bigr),\qquad
 \lim_{p\to\infty}p^6\mathcal D(p)=\frac{15}{16}.
\]
这里改进的是统一层析上包络，没有确定真实泛函的六阶系数。
\end{corollary}
\begin{proof}
由中心弱型估计~\eqref{mom:weak}，每个质量一输入的光晕面积至多
$4\alpha^{-2}$；在 $[a_0,1]$ 上改用定理~\ref{pb:volume}，并在
$[99/100,1]$ 上再减去定理~\ref{pb:sixth-gap}给出的 $(1-\alpha)^6/768$。
逐输入层析后对输入取上确界，即得所列不等式。
完整 Beta 积分为
\[
 \frac p{768}\int_0^1\alpha^{p-1}(1-\alpha)^6\,d\alpha
 =\frac{15}{16}\prod_{j=1}^6(p+j)^{-1},
\]
舍去的 $[0,99/100]$ 积分介于 $0$ 与 $(99/100)^p/768$ 之间，
故余项和极限成立。
\end{proof}


\begin{corollary}[抛物邻域对高矩包络的改进]\label{mom:parabolic-envelope}
仍记 $\mathcal B(p)$ 如推论~\ref{mom:sixth-envelope}。则对 $p>2$，
\[
 \mathfrak F_{1,\mathcal Q_c}(p)\le\mathcal B(p)-\mathcal D_*(p),\qquad
 \mathcal D_*(p)=2^{-24}p\int_{99/100}^1\alpha^{p-1}(1-\alpha)^{9/2}\,d\alpha,
\]
且
\[
 \lim_{p\to\infty}p^{9/2}\mathcal D_*(p)
 =2^{-24}\Gamma(11/2)=\frac{945\sqrt\pi}{2^{29}}.
\]
\end{corollary}
\begin{proof}
逐输入使用定理~\ref{pb:parabolic-gap}并重复前一推论的层析论证。
完整积分等于 $2^{-24}pB(p,11/2)$，截去的低端至多为 $2^{-24}(99/100)^p$。
由 Beta 积分的 Gamma 比值渐近即得极限。
这仍是统一上包络的减项，而非真实泛函的高阶系数。
\end{proof}


\begin{corollary}[高矩包络的四阶减项]\label{mom:quartic-envelope}
存在绝对 $\kappa,t_0>0$，使对 $p>2$，
\[
 \mathfrak F_{1,\mathcal Q_c}(p)\le\mathcal B(p)-\mathcal D_4(p),\qquad
 \mathcal D_4(p)=\kappa p\int_{1-t_0}^1\alpha^{p-1}(1-\alpha)^4\,d\alpha,
\]
且 $\lim_{p\to\infty}p^4\mathcal D_4(p)=24\kappa>0$。
这里可将 $t_0$ 缩小到 $t_0<1-1/\sqrt2$。
\end{corollary}
\begin{proof}
逐输入使用定理~\ref{pb:quartic-gap}，然后层析并对输入取上确界。
完整 Beta 积分为 $24\kappa\prod_{j=1}^4(p+j)^{-1}$，
截去部分至多为 $\kappa(1-t_0)^p$。
仍未据此确定真实泛函的四阶系数。
\end{proof}

\begin{corollary}[三阶亏损对高矩次项的严格改善]\label{mom:cubic-envelope}
取定理~\ref{ic:cubic-gap}中的 $\kappa,t_0$，并缩小 $t_0$ 使 $1-t_0>1/\sqrt2$。
则对 $p>2$，
\[
 \mathfrak F_{1,\mathcal Q_c}(p)\le\mathcal B(p)-\mathcal D_3(p),\qquad
 \mathcal D_3(p)=\kappa p\int_{1-t_0}^1\alpha^{p-1}(1-\alpha)^3\,d\alpha,
\]
且 $p^3\mathcal D_3(p)\to6\kappa>0$。特别地
\[
 0\le\liminf_{p\to\infty}p^3[\mathfrak F_{1,\mathcal Q_c}(p)-4p^{-2}]
 \le\limsup_{p\to\infty}p^3[\mathfrak F_{1,\mathcal Q_c}(p)-4p^{-2}]
 \le8-6\kappa<8.
\]
\end{corollary}
\begin{proof}
逐输入应用三阶亏损，再按层析公式积分并取上确界。
完整 Beta 积分为 $6\kappa\prod_{j=1}^3(p+j)^{-1}$，
截去的低端至多为 $\kappa(1-t_0)^p$。
结合推论~\ref{mom:pair-bracket}证明中的展开及固定 $Q_0$ 下界，即得次项夹逼。
这改善了真实泛函的统一上界，尚未给出其精确三阶系数。
\end{proof}

\section{分数阶消失时的双参数过渡}
\label{sec:vanishing-order}

本节讨论 $\beta\downarrow0$ 与 $\alpha\uparrow1$ 的相互作用。结论只提供已证明的
上控制和固定方体下剖面；它们不能替代全体集合上的统一双参数渐近。

设 $0<\beta<n$，$\theta=\beta/n$，且 $0<|E|<\infty$。对一个严格见证 $B$，令
\[
 s:=|E\cap B|,\qquad d:=\frac{s}{|B|},\qquad q:=\frac{s}{|E|}.
\]
则 $0<d,q\leq1$，且归一化见证值有恒等式
\[
 \frac{|B|^{\theta-1}s}{|E|^\theta}=d^{1-\theta}q^\theta.
\]

\begin{proposition}[经典光晕的统一上控制]\label{van:classical-control}
对任意三种基 $\mathcal B\in\{\mathcal Q,\mathcal R,\mathcal Q_c\}$，
\[
 H_{\beta,\mathcal B}(\alpha,E)
 \subseteq
 \left\{x:M_{0,\mathcal B}\mathbf1_E(x)>\alpha^{1/(1-\theta)}\right\}.
\]
\end{proposition}

\begin{proof}
若 $B$ 是左端任一点的严格见证，则
$d^{1-\theta}q^\theta>\alpha$。因 $q\leq1$，有 $d^{1-\theta}>\alpha$，从而
$d>\alpha^{1/(1-\theta)}$。同一 $B$ 因而见证右端严格水平集。
\end{proof}

\begin{proposition}[中心方体的过渡下剖面]\label{van:centered-slice}
对 $Q_0=[-1/2,1/2]^n$，
\[
 C_{\beta,\mathcal Q_c}(\alpha)
 \geq\bigl(1-\alpha^{1/\beta}\bigr)^n.
\]
特别地，固定 $\tau>0$ 并取 $\alpha_\beta=e^{-\beta\tau}$，则
\[
 \liminf_{\beta\downarrow0}
 C_{\beta,\mathcal Q_c}(\alpha_\beta)
 \geq(1-e^{-\tau})^n.
\]
\end{proposition}

\begin{proof}
在 $Q_0$ 内取以 $x$ 为中心、边长严格大于 $\alpha^{1/\beta}$ 的方体。
只要 $\|x\|_\infty<(1-\alpha^{1/\beta})/2$，这样的方体存在且完全包含在 $Q_0$ 内，
其分数次平均严格大于 $\alpha$。故相应中心光晕包含边长
$1-\alpha^{1/\beta}$ 的方体。令 $\alpha=\alpha_\beta$ 即得第二式。
\end{proof}

\begin{remark}[已知边界与开放问题]\label{van:open-problem}
命题~\ref{van:classical-control} 对任意输入给出经典上控制，命题~\ref{van:centered-slice}
只给出固定 $Q_0$ 的下剖面。它们显示自然过渡尺度为 $(1-\alpha)/\beta$：在
$\alpha=e^{-\beta\tau}$ 下，固定 $\tau$ 保留非零中心光晕体积。因此，不能从固定
$\beta$ 的 $\alpha\uparrow1$ 渐近推出关于 $\beta\downarrow0$ 的一致结论，也不能交换
这两个极限。仍待解决的是：在 $\tau$ 有界或 $\tau\to\infty$ 的区域，对每一种几何基
建立匹配的全体集合上下界与统一余项。
\end{remark}

\section{势函数比较与直接应用}\label{app:section}

本节给出前述几何估计的直接应用。所有结论的方向都是“势函数或解控制分数次极大函数”；因此它们产生超水平集下界和资料集中限制，并不产生新的正则性上界。

\subsection{Riesz 势的方体比较}

对 $0<\beta<n$ 和 $f\ge0$，取未归一化 Riesz 势
\[
 I_\beta f(x)=\int_{\mathbb R^n}|x-y|^{\beta-n}f(y)\,dy.
\]

\begin{theorem}[Riesz 势比较及最佳常数]\label{app:riesz}
对每个 $f\ge0$，
\[
 I_\beta f(x)\ge n^{-(n-\beta)/2}M_{\beta,\mathcal Q}f(x),
\]
以及
\[
 I_\beta f(x)\ge\left(\frac{2}{\sqrt n}\right)^{n-\beta}
 M_{\beta,\mathcal Q_c}f(x).
\]
两个常数在“对所有非负 $f$、每一点成立”的意义下最佳。
\end{theorem}
\begin{proof}
若边长为 $\ell$ 的方体 $Q$ 包含 $x$，则 $|x-y|\le\sqrt n\ell$ 对 $y\in Q$ 成立。因 $\beta-n<0$，
\[
 \int_Q|x-y|^{\beta-n}f(y)\,dy
 \ge n^{-(n-\beta)/2}\ell^{\beta-n}\int_Qf.
\]
对允许的 $Q$ 取上确界即得第一个不等式。中心方体满足 $|x-y|\le\sqrt n\ell/2$，从而得到第二个不等式。

取 $x=0$，$y_0=(1,\ldots,1)$，以及
\[
 f_\varepsilon=\varepsilon^{-n}\mathbf1_{y_0+[-\varepsilon,0]^n}.
\]
支配收敛给出 $I_\beta f_\varepsilon(0)\to n^{-(n-\beta)/2}$。任何同时含 $0$ 且与该支集相交的非中心方体边长至少 $1-\varepsilon$，故 $M_{\beta,\mathcal Q}f_\varepsilon(0)\to1$。中心方体的最小边长趋于 $2$，故 $M_{\beta,\mathcal Q_c}f_\varepsilon(0)\to2^{\beta-n}$。这两个比值分别迫使所给常数不能增大。
\end{proof}

\begin{remark}\label{app:factor}
中心与非中心 Riesz 常数之比为 $2^{n-\beta}$。它与定理~\ref{mom:centered-weak} 中弱型常数右端的同一数值因子来自不同证明：前者来自核在中心方体上的直径，后者来自有限选择的重叠数。
\end{remark}

\subsection{Wolff 势的方体比较}

设 $1<p<n$，$\mu$ 为非负 Radon 测度，定义
\[
 W_{1,p}^{\mu}(x)=\int_0^\infty
 \left(\frac{\mu(B(x,r))}{r^{n-p}}\right)^{1/(p-1)}\frac{dr}{r}.
\]

\begin{theorem}[Wolff 势比较及最佳常数]\label{app:wolff}
令 $\gamma=(n-p)/(p-1)$。对每个 $x$，
\[
 W_{1,p}^{\mu}(x)\ge
 \frac{p-1}{n-p}\,n^{-\gamma/2}
 \bigl(M_{p,\mathcal Q}\mu(x)\bigr)^{1/(p-1)},
\]
并且
\[
 W_{1,p}^{\mu}(x)\ge
 2^\gamma\frac{p-1}{n-p}\,n^{-\gamma/2}
 \bigl(M_{p,\mathcal Q_c}\mu(x)\bigr)^{1/(p-1)},
\]
其中
\[
 M_{p,\mathcal B}\mu(x)=
 \sup_{Q\in\mathcal B(x)}|Q|^{p/n-1}\mu(Q).
\]
两个常数均最佳。
\end{theorem}
\begin{proof}
若 $Q\ni x$ 的边长为 $\ell$，令 $R=\sqrt n\ell$。在 Wolff 势中保留 $[R,\infty)$ 的积分，得到
\[
 W_{1,p}^{\mu}(x)\ge
 \mu(Q)^{1/(p-1)}\int_R^\infty r^{-\gamma}\frac{dr}{r}
 =\frac1\gamma n^{-\gamma/2}
 \bigl(\ell^{p-n}\mu(Q)\bigr)^{1/(p-1)}.
\]
取上确界即得第一个结论；中心情形中以 $R=\sqrt n\ell/2$ 代替即可。

若 $\mu=\delta_y$，则 $W_{1,p}^{\delta_y}(x)=\gamma^{-1}|x-y|^{-\gamma}$。选择 $x-y$ 的所有坐标绝对值相等时，$|x-y|=\sqrt n\|x-y\|_\infty$，而非中心和中心最小见证方体的边长分别为 $\|x-y\|_\infty$ 和 $2\|x-y\|_\infty$。直接代入表明两个常数均不能提高。
\end{proof}

\subsection{两个直接后果}

\begin{corollary}[分数 Laplace 资料的非退化区域]\label{app:fractional-laplace}
设 $0<\beta<\min\{2,n\}$，$E$ 是有界可测集，并令
\[
 u(x)=c_{n,\beta}\,I_\beta\mathbf1_E(x),\qquad
 c_{n,\beta}=\frac{\Gamma((n-\beta)/2)}{2^\beta\pi^{n/2}\Gamma(\beta/2)}.
\]
采用 Fourier 约定 $\widehat f(\xi)=\int_{\mathbb R^n}e^{-ix\cdot\xi}f(x)\,dx$，并以
\[
 \widehat{(-\Delta)^{\beta/2}f}(\xi)=|\xi|^\beta\widehat f(\xi)
\]
定义分数 Laplace 算子。此约定下，上述常数满足
$(-\Delta)^{\beta/2}u=\mathbf1_E$ 于缓增分布意义；该归一化及等价定义见 Kwa\'snicki \cite{Kwasnicki2017}。若 $E$ 是轴平行方体，则每个 $0<\alpha<1$ 满足
\[
 \bigl|\{u>c_{n,\beta}n^{-(n-\beta)/2}\alpha |E|^{\beta/n}\}\bigr|
 \ge (2\alpha^{-1/(n-\beta)}-1)^n|E|,
\]
并且
\[
 \bigl|\{u>c_{n,\beta}(2/\sqrt n)^{n-\beta}\alpha |E|^{\beta/n}\}\bigr|
 \ge \max\!\left\{(1-\alpha^{1/\beta})^n,
 (\alpha^{-1/(n-\beta)}-1)^n\right\}|E|.
\]
\end{corollary}
\begin{proof}
定理~\ref{app:riesz} 把相应方体光晕包含在 $u$ 的超水平集内。再使用非中心和中心方体的实际测试方体包含，即得两个不等式。
\end{proof}

\begin{corollary}[点态有界性排除资料集中]\label{app:concentration}
若 $f\ge0$ 且 $c_{n,\beta}I_\beta f(x)\le U<\infty$，则每个含 $x$ 的轴平行方体 $Q$ 满足
\[
 \int_Q f\le \frac{U}{c_{n,\beta}n^{-(n-\beta)/2}}\ell(Q)^{n-\beta}.
\]
对以 $x$ 为中心的方体，$n^{-(n-\beta)/2}$ 可改为较大的 $(2/\sqrt n)^{n-\beta}$。
\end{corollary}
\begin{proof}
由定理~\ref{app:riesz} 控制 $M_{\beta,\mathcal Q}f(x)$，再取单个方体作为该上确界的候选即可。
\end{proof}

\begin{corollary}[拟线性资料的集中必要条件]\label{app:plaplace}
设 $1<p<n$。假设非负 $p$-超调和函数 $u$ 已按 $\inf_{\mathbb R^n}u=0$ 归一化，Riesz 测度为 $\mu$，且在所考虑点满足可用的下势估计
\[
 u(x)\ge c_0W_{1,p}^{\mu}(x).
\]
若 $u(x)\le U$，则每个 $Q\ni x$ 有
\[
 \mu(Q)\le
 \left(\frac{U}{c_0\gamma^{-1}n^{-\gamma/2}}\right)^{p-1}
 \ell(Q)^{n-p}.
\]
中心方体中相应常数可提高 $2^\gamma$ 倍。
\end{corollary}
\begin{proof}
由定理~\ref{app:wolff} 解出 $M_{p,\mathcal Q}\mu(x)$，再取 $Q$ 为上确界的一个候选。这里的结论以所述下势估计和 Wolff 势在该点有限为前提；它并不对任意有限测度自动成立。
\end{proof}

\begin{remark}[应用边界]\label{app:boundary}
上述推论只给出了“解大于势下界”导出的必要条件和超水平集下界。Schauder、Calder\'on--Zygmund、梯度高可积性或 good--$\lambda$ 正则性仍需要反向分布控制、局部比较或容量估计；它们不能由这里的单向比较推出。有关分数极大分布函数和正则性接口的已有方法见 Nguyen--Tran \cite{NguyenTran2021}。
\end{remark}

\subsection{限制型范数、反向扩张与后续问题}

定理~\ref{gl:main}的全体中心光晕首项确定了限制型高矩的精确首项；进一步的统一估计还关系到近临界
$L^{r_p}\to L^p$ 示性资料常数。Bañuelos--Osękowski 的两权问题
\cite{BanuelosOsekowski2025}研究树状概率空间上固定指数的全函数范数；与本文可比较的
后续目标是欧氏示性资料的最优序列及其归一化形状。另一个方向由限制反向弱型
\cite{Lerner2026}提示：在非中心模型中识别对迭代封闭的近似基元素资料类，
证明相应的分数次反扩张与稳定性，再考察它如何进入局部分布估计。

本文已求得固定方体和任意固定锚定凸体的切向极限，并给出二维固定方体的高水平精确公式。
对全体可测输入，中心方体首项已经匹配；更精细的常数与一般几何基仍需研究。
二维半阶中，定理~\ref{pb:volume}将体积三阶剩余量夹在 $[2,10/3]$，
推论~\ref{mom:pair-bracket}将高矩次阶剩余量夹在 $[0,8]$；
定理~\ref{pb:strict-gap}已证明每个固定高阈值的两点上包络具有严格间隙；其端点衰减阶、极限存在性及上下界匹配仍未确定。具体问题包括：
\begin{enumerate}
\item 求高维非中心方体与强矩形的精确首项；确定中心方体的全体下一系数，并刻画可能并列的近极值构型。
\item 研究凸体基上对全体输入的匹配常数和仿射极值体。固定输入 $E=K$ 的联合切向体已由定理~\ref{cvx:full-tangent-limit}确定，对称半阶模型的椭球极值已由推论~\ref{half:ellipsoid-extremum}确定。
\item 在 $\beta\downarrow0$ 与 $\alpha\uparrow1$ 的耦合尺度中证明统一渐近；在二维精细公式的临界分数阶附近建立一致过渡。
\item 将近重合几何扩展到允许旋转、无界偏心率的几何族，或齐次群与 Ahlfors 正则空间。
\item 建立内禀一般测度模型的非中心与强矩形覆盖理论，研究度量球上一致弱型常数及极值结构。
\item 确定中心限制型高矩的下一系数及非中心精确首项，建立分数次反限制弱型迭代及其近极值理论。
\item 确定一般多输入几何平均模型的有限阈值常数，判断近极值资料是否渐近对角。
\item 构造带可控示性分解余项的反向 good--$\lambda$ 或容量估计，并在具体测度资料方程中检验定量正则性后果。
\end{enumerate}

\bibliographystyle{unsrtnat}
\bibliography{references}
\end{document}
